\documentclass{article}
\usepackage[colorlinks=true, linkcolor=blue, citecolor=black, urlcolor=black]{hyperref}
\usepackage{amsmath,amsthm,amssymb,dsfont}
\usepackage{thmtools}
\usepackage{algorithm,algpseudocodex}
\usepackage{cleveref}
\usepackage[section]{placeins}
%\usepackage{draftwatermark}
%\SetWatermarkScale{8}
%\SetWatermarkText{\normalfont{DRAFT}}

\newcommand{\PairDecomp}{\mathcal{W}}

\title{TSP and WSP Thoughts and Ideas}
\author{Addison Hanrattie}

\newtheorem{thm}{Theorem}[section]
\newtheorem{lem}[thm]{Lemma}
\newtheorem{corollary}{Corollary}[thm]
\newtheorem*{rem*}{Remark}
\theoremstyle{definition}
\newtheorem{defn}[thm]{Definition}
\crefname{thm}{Theorem}{Theorems}
\crefname{lem}{Lemma}{Lemmas}
\crefname{corollary}{Corollary}{Corollaries}


\newcommand{\BigLength}{\mathcal{L}}
\newcommand{\Cover}{\mathcal{C}}

\newcommand{\dist}{\delta}
%\DeclareMathOperator{\dist}{dist}

\DeclareMathOperator{\diam}{diam}

\begin{document}
\maketitle

\tableofcontents

\section{Work Area}
The main theorem I am proving here is Corollary~\ref{cor:wsp_tsp} and Theorem~\ref{thm:bin_wspd_tour_structure}.

\section{Definitions}
\subsection{Background Definitions}
\begin{defn}
    A \textbf{Pair Decomposition} of a point set $P$ is a set of pairs
    $$\PairDecomp = \left\{ \{A_1, B_1\}, \ldots, \{A_k,B_k\} \right\}$$
    such that (I) $\forall i; A_i,B_i \subset P$, (II) $\forall i; A_i \cap B_i = \emptyset$, and (III) $\bigcup_{i=1}^k A_i \otimes B_i = \binom{P}{2} = P \otimes P$.
\end{defn}

\begin{defn}
    $\dist(Q,R) = \min_{q \in Q, r \in R} |qr|$
\end{defn}

\begin{defn}
    The pair $\{Q,R\}$ (where $Q \cap R = \emptyset$) is \textbf{$s$-well-separated} if
    $$\max\{\diam(Q), \diam(R)\} \leq \frac{1}{s} \cdot \dist(Q,R)$$
\end{defn}

\begin{corollary}
    Given two sets of points $A$ and $B$, there exists an $s$ such that $A$ and $B$ are $s$-well-separated if and only if $\dist(A,B) > 0$.
\end{corollary}

\subsection{My definitions}

\begin{defn}
    Let $P = A \cup B$ (with $A \cap B = \emptyset$). We say that $A$ and $B$ are \textbf{TSP-Separated} if there exists an optimal TSP tour on $P$ that visits all points in $A$ before any point in $B$ (or vice versa). 
    Furthermore, Let $P = A \cup B \cup C$ (with all three pairwise disjoint). We say that $A$, $B$, and $C$ are \textbf{TSP-Separated} if 
    there exists an optimal TSP tour on $P$ that has exactly 3 or 4 edges connecting the three sets. 
\end{defn}

\begin{defn}
    Let $\{(A_1, B_1), \dots, (A_n, B_n)\} = \PairDecomp$ be a Well Separated Pair Decomposition of the point set $P$.
    We say that a given pair $(A_i, B_i) \in \PairDecomp$ is a \textbf{$k$-binary WSP} if either $A_i \cup B_i = P$ (which we call $1$-binary WSP) or there exists another pair $(A_j, B_j) \in \PairDecomp$ such that $A_i \cup B_i$ is equal to either $A_j$ or $B_j$ and furthermore the pair $(A_j, B_j)$ is a $(k-1)$-binary WSP.
\end{defn}

\begin{defn}
    A Well Separated Pair Decomposition $\PairDecomp$ of the point set $P$ is a \textbf{$k$-binary WSPD} if the union of all $k$-binary WSPs in $\PairDecomp$ is equal to $P$.
\end{defn}

\begin{defn}
    A Well Separated Pair Decomposition $\PairDecomp$ of the point set $P$ is a \textbf{full binary WSPD} if $\forall (A_i, B_i) \in \PairDecomp$ the pair is a $k$-binary WSP for some $k$.
\end{defn}

\section{Background Theorems}
\begin{thm}
Let $A$ and $B$ be two sets of points that are $s$-well-separated, let $p,p' \in A$, and let $q,q' \in B$. 
Then the following inequalities hold:
\begin{align*}
    |pq| &\leq (1 + 1/s) \cdot |pq'| \\
    |pq| &\leq (1 + 2/s) \cdot |p'q'| \\
    |pp'| &\leq (1/s) \cdot |pq| \\
\end{align*}
\end{thm}

\begin{proof}
    Let $A$ and $B$ be two sets of points that are $s$-well-separated, let $p,p' \in A$, and let $q,q' \in B$.

    Note that $|qq'| \leq \diam(B) \leq \frac{1}{s} \cdot \dist(A,B) \leq \frac{1}{s} |pq'|$ by the definition of well-separated pairs.
    The first result then follows from the next analysis:
    \begin{align*}
        |pq| &\leq |pq'| + |q'q| \\
        &\leq |pq'| + \frac{1}{s} |pq'| \\
        &= (1 + 1/s) \cdot |pq'|
    \end{align*}

    For the second result we note that $|pp'| \leq \diam(A) \leq \frac{1}{s} \cdot \dist(A,B) \leq \frac{1}{s} |p'q'|$ by the definition of well-separated pairs.
    The second result then follows from the next analysis:
    \begin{align*}
        |pq| &\leq |pp'| + |p'q'| + |q'q| \\
        &\leq |pp'| + |p'q'| + \frac{1}{s} |p'q'| \\
        &\leq \frac{1}{s} |p'q'| + |p'q'| + \frac{1}{s} |p'q'| \\
        &= (1 + 2/s) \cdot |p'q'|
    \end{align*}
    
    Finally for the third result we note that the following is true $|pp'| \leq \diam(A) \leq \frac{1}{s} \cdot \dist(A,B) \leq \frac{1}{s} |pq|$ by the definition of well-separated pairs. The result is immediately aparent that $|pp'| \leq (1/s) \cdot |pq|$.
\end{proof}

\section{My Theorems}

\begin{lem}\label[lem]{lem:wsp_subset}
    Let $A$ and $B$ be two sets of points that are $s$-well-separated.
    Then $A' \subseteq A$ and $B' \subseteq B$ (where $|A'|, |B'| \geq 1$) are also $s$-well-separated.
\end{lem}
\begin{proof}
    Let $A$ and $B$ be two sets of points that are $s$-well-separated and let $A' \subseteq A$ and $B' \subseteq B$ (where $|A'|, |B'| \geq 1$).
    We can then note that $\dist(A', B') \geq \dist(A,B)$ due to the triangle inequality.
    We can also note that $\diam(A') \leq \diam(A)$ and $\diam(B') \leq \diam(B)$ since $A' \subseteq A$ and $B' \subseteq B$.
    Thus we have the following:
    $$\max\{\diam(A'), \diam(B')\} \leq \max\{\diam(A), \diam(B)\} \leq \frac{1}{s} \cdot \dist(A,B) \leq \frac{1}{s} \cdot \dist(A', B')$$
    as desired.
\end{proof}

\subsection{Hamiltonian Endpoint Distance Bound}

\begin{thm}
Within a ball of diameter $d$ within a metric space the difference in shortest hamiltonian paths between any two points is at most $2d$.
\end{thm}
\begin{proof}
    Fix one pair of endpoints $(a,b)$ and let $P_{a,b}$ be the optimal Hamiltonian path from $a$ to $b$ with length denoted $L(a,b)$.
    Next consider another pair of endpoints $(c,d)$ we can then build a hamiltonian path from $c$ to $d$ using $P_{a,b}$.
    \begin{enumerate}
        \item Start at $c$
        \item Travel from $c$ to $a$ for a cost of $\dist(c,a) \leq d$
        \item Follow the path $P_{a,b}$ from $a$ to $b$ for a cost of $c \leq L(a,b)$. (shortcutting $c,d$)
        \item Travel from $b$ to $d$ for a cost of $\dist(b,d) \leq d$
    \end{enumerate}
    %WLOG assume that $c$ occurs before $d$ on the path $P_{a,b}$. 
    %Then let $x$ be the point which precedes $c$ on the path and let $x'$ be the point that immediately follows. 
    %In the same way define $y$ and $y'$ for point $d$.
    %Since the triangle inequality must be obeyed we know 
    Since the space obeys the triangle inequality we know that during step 3 we can only decrease the cost by shortcutting. Thus the total cost of this path is at most $L(a,b) + 2d$. Furthermore we know that this cost must ceil whatever the optimal hamiltonian path from $c$ to $d$ is. Thus we have shown that $L(c,d) \leq L(a,b) + 2d$. By symmetry of the argument we can also show that $L(a,b) \leq L(c,d) + 2d$. Combining these two inequalities we have that $\forall a,b,c,d;|L(a,b) - L(c,d)| \leq 2d$ as desired. 
\end{proof}
I haven't been able to prove it either way yet, but I believe the bound can be tightened for this theorem. 


\begin{thm}
    Given a set of points $P$ within a ball of diameter $d$ within a metric space and starting point $x,y \in P$. Let $L$ be the length of some optimal hamiltonian path $H$ between $x$ and $y$ over $P$. 
    Then the length of the optimal hamiltonian path after removing $k$ intermediate points is at least $L - d(k+b)$ where $b$ is the number of non-contiguous blocks of points removed (under $H$).
\end{thm}
\begin{proof}
    Let $Q \subseteq P \setminus \{x,y\}$ be the set of $k$ intermediate points that we remove from $P$.
    We then partition $Q$ into $b$ blocks of consecutive vertices such that $B_j = \{v_{j,1}, \dots, v_{j,m_j}\}$ is the $j$-th block of points in $Q$ and furthermore there are no edges in $H$ between any two blocks of points in $Q$. 
    Thus we will have $\sum_j m_j = k$.

    We define $H'$ to be an optimal hamiltonian path between $x$ and $y$ over $P \setminus Q$ and to have length $L'$.

    We can then build a new hamiltonian path $\widehat H$ by inserting the blocks of points in $Q$ back into $H'$ in their same order as they appear in $H$.
    For a block $B_j$ we select any edge $(u_1, u_2)$ in the current version of $\widehat H$ and replace it with 
    $(u_1, v_{j,1}), (v_{j,1}, v_{j,2}), \dots, (v_{j,m_j}, u_2)$.
    The increase in length will be
    $$\dist(u_1, v_{j,1}) + \sum_{i=1}^{m_j-1} \dist(v_{j,i}, v_{j,i+1}) + \dist(v_{j,m_j}, u_2) - \dist(u_1, u_2)$$
    
    Since all distances are at most $d$, the increase is at most $(m_j + 1)d$.
    Then summing over all $b$ blocks we have the total increase at most $\sum_{j=1}^b (m_j + 1)d = d(k+b)$.

    Then since $H$ is optimal we have that $L \leq L' + d(k+b)$ which implies that $L' \geq L - d(k+b)$ as desired.
\end{proof}

\begin{corollary}
    Given a set of points $P$ within a ball of diameter $d$ within a metric space and starting point $x,y \in P$. 
    Let $L$ be the length of some optimal hamiltonian path $H$ between $x$ and $y$ over $P$. 
    Then the length of the optimal hamiltonian path after removing $k$ intermediate points is at least $L - 2dk$.
\end{corollary}

\subsection{Well-Separated Sets are Path/Tour Separated}
\subsubsection{Single WSP Single Path Analysis}

\input{../MSThesis/theorems/wsp_path_ab.tex}

\input{../MSThesis/theorems/wsp_path_aa.tex}

\input{../MSThesis/theorems/wsp_tsp.tex}

\input{../MSThesis/theorems/wsp_tsp_tightness.tex}

\subsubsection{Single WSP Endpoint Pair Multi Path Analysis}

\begin{lem}\label[lem]{lem:wsp_twopath_fixed_aa_aa}
    Suppose we have two sets of points $A$ and $B$ that are $s$-well-separated where $s \geq \frac{3}{2}$.
    Then the multiple optimal hamiltonian paths problem where one path had fixed endpoints $a', a'' \in A$ and the other path had fixed endpoints $a''', a'''' \in A$ has a solution where one of the path visits all points in $B$ consecutively and the other path has no edges connecting $A$ to $B$.
    Furthermore if $s > \frac{3}{2}$ then all optimal solutions must satisfy this property.
\end{lem}
\begin{proof} % gemini checked for major flaws
    Consider some optimal solution to the problem.
    Suppose that both paths have edges connecting $A$ to $B$.
    Then there must be two edges of the form $(a^*_1, b^*_1)$ and $(b^*_2, a^*_2)$ in the first set and two edges of the form $(a^\dagger_1, b^\dagger_1)$ and $(b^\dagger_2, a^\dagger_2)$ in the second set where $a^*_1, a^*_2, a^\dagger_1, a^\dagger_2 \in A$ and $b^*_1, b^*_2, b^\dagger_1, b^\dagger_2 \in B$.
    Thus these paths have the following partial structure:
    $$a' \leadsto a^*_1 \to b^*_1 \leadsto b^*_2 \to a^*_2 \leadsto a''
    \quad\quad
    a''' \leadsto a^\dagger_1 \to b^\dagger_1 \leadsto b^\dagger_2 \to a^\dagger_2 \leadsto a''''$$

    Since the solution is optimal we know that they satisfy the 3-optimality condition implying the following:
    $|a^*_1 b^*_1| + |b^*_2 a^*_2| + |a^\dagger_1 b^\dagger_1| \leq |a^*_1 a^*_2| + |a^\dagger_1 b^*_1| + |b^*_2 b^\dagger_1|$.
    This is to say that the following paths must have a combined length at least as long as the original paths:
    $$a' \leadsto a^*_1 \to a^*_2 \leadsto a'' \quad\quad a''' \leadsto a^\dagger_1 \to b^*_1 \leadsto b^*_2 \to b^\dagger_1 \leadsto b^\dagger_2 \to a^\dagger_2 \leadsto a''''$$

    We can define $L$ to be the combined length of the original paths and $L'$ to be the combined length of the newly constructed paths.
    We can then note the following about the change in length between $L$ and $L'$:
    $L - L' = (|a^*_1 b^*_1| + |b^*_2 a^*_2| + |a^\dagger_1 b^\dagger_1|) - (|a^*_1 a^*_2| + |a^\dagger_1 b^*_1| + |b^*_2 b^\dagger_1|)$.

    Define $\omega = \dist(A,B)$.
    Using the definitions of well separation we have:
    \[
        |b^*_2 a^*_2| \geq \omega \quad\quad
        |a^\dagger_1 b^\dagger_1| \geq \omega \quad\quad
        |a^*_1 a^*_2| \leq \frac{1}{s} \omega \quad\quad
        |a^\dagger_1 a^*_1| \leq \frac{1}{s} \omega \quad\quad
        |b^*_2 b^\dagger_1| \leq \frac{1}{s} \omega
    \]

    The following can then be determined based on these inequalities and the properties of a metric space:
    \begin{align*}
        L - L' &= (|a^*_1 b^*_1| + |b^*_2 a^*_2| + |a^\dagger_1 b^\dagger_1|) - (|a^*_1 a^*_2| + |a^\dagger_1 b^*_1| + |b^*_2 b^\dagger_1|)\\
        &\geq (|a^*_1 b^*_1| + |b^*_2 a^*_2| + |a^\dagger_1 b^\dagger_1|) - (|a^*_1 a^*_2| + |a^\dagger_1 a^*_1| + |a^*_1 b^*_1| + |b^*_2 b^\dagger_1|) \\
        &= (|b^*_2 a^*_2| + |a^\dagger_1 b^\dagger_1|) - (|a^*_1 a^*_2| + |a^\dagger_1 a^*_1| + |b^*_2 b^\dagger_1|) \\
        &\geq 2\omega - (|a^*_1 a^*_2| + |a^\dagger_1 a^*_1| + |b^*_2 b^\dagger_1|) \\
        &\geq 2\omega - \frac{3}{s} \omega \\
    \end{align*}

    If $s > \frac{3}{2}$ then we have that $L - L' > 0$ which is a contradiction since if the original path were optimal then we would have $L - L' \leq 0$.
    Therefore no optimal path can have both paths visiting both $A$ and $B$ if $s > \frac{3}{2}$.
    If $s = \frac{3}{2}$ then we have that $L - L' \geq 0$ which combined with $L - L' \leq 0$ (by optimality) implies that $L - L' = 0$.
    Therefore the newly constructed paths must also be optimal solutions to the problem.
    Since the number of crossings between $A$ and $B$ has strictly decreased we can apply this logic repeatedly to construct an optimal solution where only a single path has edges connecting $A$ to $B$ and the other path has no edges connecting $A$ to $B$ as desired.

    We can then consider a case where only a single path has edges connecting $A$ to $B$.
    We can then mark $A'$ to be the set of points in $A$ that are visited by this path and $B'$ to be the set of points in $B$ that are visited by this path.
    Then we can note that $A'$ and $B'$ are still $s$-well-separated by \Cref{lem:wsp_subset}.
    We can then apply \Cref{thm:wsp_path_aa} to construct a solution where this path visits all points in $B$ consecutively if it was not already the case.
\end{proof}

\begin{lem}\label[lem]{lem:wsp_twopath_fixed_aa_ab}
    Suppose we have two sets of points $A$ and $B$ that are $s$-well-separated where $s \geq \frac{3}{2}$.
    Then the multiple optimal hamiltonian paths problem where one path had fixed endpoints $a', a'' \in A$ and the other path had fixed endpoints $a''' \in A$ and $b' \in B$ has a solution where the first path visits only points in $A$ and the second path visits all points in $B$ consecutively.
    Furthermore if $s > \frac{3}{2}$ then all optimal solutions must satisfy this property.
\end{lem}
\begin{proof} % gemini checked for major flaws
    Begin by considering some optimal solution to the problem where the first path visits some points in $B$.
    Then we can note the following edges must exist in the first path $(a^*_1, b^*_1)$ and $(b^*_2, a^*_2)$ where $a^*_1, a^*_2 \in A$ and $b^*_1, b^*_2 \in B$ and in the second path there must be an edge of the form $(a^\dagger, b^\dagger)$ where $a^\dagger \in A$ and $b^\dagger \in B$.
    Thus these paths have the following partial structure:
    $$a' \leadsto a^*_1 \to b^*_1 \leadsto b^*_2 \to a^*_2 \leadsto a''
    \quad\quad
    a''' \leadsto a^\dagger \to b^\dagger \leadsto b'$$

    Since the solution is optimal we know that they satisfy the 3-optimality condition implying the following:
    $|a^*_1 b^*_1| + |b^*_2 a^*_2| + |a^\dagger b^\dagger| \leq |a^*_1 a^*_2| + |a^\dagger b^*_1| + |b^*_2 b^\dagger|$.
    This is to say that the following paths must have a combined length at least as long as the original paths:
    $$a' \leadsto a^*_1 \to a^*_2 \leadsto a'' \quad\quad a''' \leadsto a^\dagger \to b^*_1 \leadsto b^*_2 \to b^\dagger \leadsto b'$$

    We can define $L$ to be the combined length of the original paths and $L'$ to be the combined length of the newly constructed paths.
    We can then note the following about the change in length between $L$ and $L'$:
    $L - L' = (|a^*_1 b^*_1| + |b^*_2 a^*_2| + |a^\dagger b^\dagger|) - (|a^*_1 a^*_2| + |a^\dagger b^*_1| + |b^*_2 b^\dagger|) \leq 0$ by optimality.

    Define $\omega = \dist(A,B)$.
    Using the definitions of well separation and rules of a metric space we have:
    \[
        |a^*_1 b^*_1| \geq \omega \quad\quad
        |b^*_2 a^*_2| \geq \omega \quad\quad
        |a^*_1 a^*_2| \leq \frac{1}{s} \omega \quad\quad
        |b^\dagger b^*_1| \leq \frac{1}{s} \omega \quad\quad
        |b^*_2 b^\dagger| \leq \frac{1}{s} \omega \quad\quad
    \]

    The following can then be determined based on these inequalities and the properties of a metric space:
    \begin{align*}
        L - L' &= (|a^*_1 b^*_1| + |b^*_2 a^*_2| + |a^\dagger b^\dagger|) - (|a^*_1 a^*_2| + |a^\dagger b^*_1| + |b^*_2 b^\dagger|) \\
        &\geq (|a^*_1 b^*_1| + |b^*_2 a^*_2| + |a^\dagger b^\dagger|) - (|a^*_1 a^*_2| + |a^\dagger b^\dagger| + |b^\dagger b^*_1| + |b^*_2 b^\dagger|) \\
        &= (|a^*_1 b^*_1| + |b^*_2 a^*_2|) - (|a^*_1 a^*_2| + |b^\dagger b^*_1| + |b^*_2 b^\dagger|) \\
        &\geq 2\omega - (|a^*_1 a^*_2| + |b^\dagger b^*_1| + |b^*_2 b^\dagger|) \\
        &\geq 2\omega - \frac{3}{s} \omega \\
    \end{align*}

    If $s > \frac{3}{2}$ then we have that $L - L' > 0$ which is a contradiction since if the original path were optimal then we would have $L - L' \leq 0$.
    Therefore no optimal path can have both paths visiting both $A$ and $B$ if $s > \frac{3}{2}$.
    If $s = \frac{3}{2}$ then we have that $L - L' \geq 0$ which combined with $L - L' \leq 0$ (by optimality) implies that $L - L' = 0$.
    Therefore the newly constructed paths must also be optimal solutions to the problem.
    Since the number of crossings between $A$ and $B$ has strictly decreased we can apply this logic repeatedly to construct an optimal solution where only the second path has edges connecting $A$ to $B$.

    We can then consider a case where only the second path has edges connecting $A$ to $B$.
    We can then mark $A'$ to be the set of points in $A$ that are visited by this path and $B'$ to be the set of points in $B$ that are visited by this path.
    Then we can note that $A'$ and $B'$ are still $s$-well-separated by \Cref{lem:wsp_subset}.
    We can then apply \Cref{thm:wsp_path_ab} to construct a solution where this path visits all points in $B$ consecutively if it was not already the case.
\end{proof}

\begin{lem}\label[lem]{lem:wsp_twopath_fixed_ab_ab}
    Suppose we have two sets of points $A$ and $B$ that are $s$-well-separated where $s \geq 1$.
    Then the multiple optimal hamiltonian paths problem where one path had fixed endpoints $a' \in A$ and $b' \in B$ and the other path had fixed endpoints $a'' \in A$ and $b'' \in B$ has a solution where each path only has a single edge connecting $A$ to $B$ and furthermore if $s > 1$ then all optimal solutions must satisfy this property.
\end{lem}
\begin{proof} % gemini checked for major flaws
    Consider some optimal solution to the problem.
    We can define $A'$ to be the set of points in $A$ that are visited by the first path and $B'$ to be the set of points in $B$ that are visited by the first path.
    We can also define $A''$ to be the set of points in $A$ that are visited by the second path and $B''$ to be the set of points in $B$ that are visited by the second path.
    Then we can note that $A'$ and $B'$ are still $s$-well-separated and $A''$ and $B''$ are still $s$-well-separated by \Cref{lem:wsp_subset}.
    We can then apply \Cref{thm:wsp_path_ab} to both of the paths which states that if $s > 1$ then the two paths could not have more than one edge connecting $A$ to $B$ and if $s = 1$ then there exists an optimal solution for each of the subproblems where they only cross $A$ and $B$ once.
\end{proof}

\begin{lem}\label[lem]{lem:wsp_twopath_fixed_aa_bb}
    Suppose we have two sets of points $A$ and $B$ that are $s$-well-separated where $s \geq \frac{3}{2}$.
    Then the multiple optimal hamiltonian paths problem where one path had fixed endpoints $a',a'' \in A$ and the other path had fixed endpoints $b',b'' \in B$ has a solution where neither path has an edge connecting $A$ to $B$ and furthermore if $s > \frac{3}{2}$ then all optimal solutions must satisfy this property.
\end{lem}
\begin{proof} % gemini checked for major flaws
    Assume that for some optimal solution it has at least one edge connecting $A$ to $B$.
    We will define $P_A$ to be the first path that connects $a'$ to $a''$ and $P_B$ to be the second path that connects $b'$ to $b''$.
    WLOG assume that $P_A$ has an edge connecting $A$ to $B$ (if not then flip the names of $A$ and $B$)
    
    \textbf{Case 1:} $P_B$ also traverses part of $A$.
    Therefore we have the following two edges in $P_A$: $(a^*_1, b^*_1), (b^*_2, a^*_2)$ and the following edges in $P_B$: $(b^\dagger_1, a^\dagger_1), (a^\dagger_2, b^\dagger_2)$ all ordered as they appear in their respective paths.
    This gives the following partial structure of the paths:
    $$a'_s \leadsto a^*_1 \to b^*_1 \leadsto b^*_2 \to a^*_2 \leadsto a'_e
    \quad\quad
    b'_s \leadsto b^\dagger_1 \to a^\dagger_1 \leadsto a^\dagger_2 \to b^\dagger_2 \leadsto b'_e$$

    Since the paths are optimal we know that they satisfy the 4-optimality condition implying the following:
    $|a^*_1 b^*_1| + |b^*_2 a^*_2| + |b^\dagger_1 a^\dagger_1| + |a^\dagger_2 b^\dagger_2| \leq |a^*_1 a^\dagger_1| + |a^\dagger_2 a^*_2| + |b^\dagger_1 b^*_1| + |b^*_2 b^\dagger_2|$.
    This is to say that the following paths must have a combined length at least as long as the original paths:
    $$a'_s \leadsto a^*_1 \to a^\dagger_1 \leadsto a^\dagger_2 \to a^*_2 \leadsto a'_e
    \quad\quad
    b'_s \leadsto b^\dagger_1 \to b^*_1 \leadsto b^*_2 \to b^\dagger_2 \leadsto b'_e$$

    Since $A$ and $B$ are $s$-well-separated we know that $\forall i,j,k; |a_i a_j| \leq \frac{1}{s} |a_i b_k|$ (and vice versa for $B$).
    From this we can conclude the following:
    \begin{align*}
        |a^*_1 b^*_1| + |b^*_2 a^*_2| + |b^\dagger_1 a^\dagger_1| + |a^\dagger_2 b^\dagger_2| &\leq |a^*_1 a^\dagger_1| + |a^\dagger_2 a^*_2| + |b^\dagger_1 b^*_1| + |b^*_2 b^\dagger_2| \\
        &\leq \frac{1}{s} |a^*_1 b^*_1| + \frac{1}{s} |b^*_2 a^*_2| + \frac{1}{s} |b^\dagger_1 a^\dagger_1| + \frac{1}{s} |a^\dagger_2 b^\dagger_2| \\
        &= \frac{1}{s} (|a^*_1 b^*_1| + |b^*_2 a^*_2| + |b^\dagger_1 a^\dagger_1| + |a^\dagger_2 b^\dagger_2|)
    \end{align*}
    This implies that $1 \leq \frac{1}{s}$ which in turn implies that $s \leq 1$.
    This however is a contradiction since we assumed that $s \geq \frac{3}{2}$.

    \textbf{Case 2:} $P_B$ does not traverse any part of $A$.
    Then we must have two edges in $P_A$: $(a^*_1, b^*_1), (b^*_2, a^*_2)$ and one edge in $P_B$: $(b^\dagger_1, b^\dagger_2)$ all ordered as they appear in their respective paths.
    Furthermore it follows from \Cref{thm:wsp_path_aa} that these must be the only two edges between $A$ and $B$.
    This gives the following partial structure of the paths:
    $$a'_s \leadsto a^*_1 \to b^*_1 \leadsto b^*_2 \to a^*_2 \leadsto a'_e
    \quad\quad
    b'_s \leadsto b^\dagger_1 \to b^\dagger_2 \leadsto b'_e$$

    Since the paths are optimal we know that they satisfy the 3-optimality condition implying the following:
    $|a^*_1 b^*_1| + |b^*_2 a^*_2| + |b^\dagger_1 b^\dagger_2| \leq |a^*_1 a^*_2| + |b^\dagger_1 b^*_1| + |b^*_2 b^\dagger_2|$
    This is to say that the following paths must have a combined length at least as long as the original paths:
    $$a'_s \leadsto a^*_1 \to a^*_2 \leadsto a'_e
    \quad\quad
    b'_s \leadsto b^\dagger_1 \to b^*_1 \leadsto b^*_2 \to b^\dagger_2 \leadsto b'_e$$

    We can define $L$ to be the total length of the two paths as is. 
    We can then build a new solution by shortcutting the path by directly connecting $a^*_1$ to $a^*_2$ within $P_A$ to form $P'_A$ and we can break the edge $(b^\dagger_1, b^\dagger_2)$ in $P_B$ and then connect $b^\dagger_1$ to $b^*_1$ and $b^*_2$ to $b^\dagger_2$ to form $P'_B$.
    Let $L'$ be the total length of the two paths after this modification.
    We can then note the following about the reduction in length from this modification:
    $L - L' = (|a^*_1 b^*_1| + |b^*_2 a^*_2| + |b^\dagger_1 b^\dagger_2|) - (|a^*_1 a^*_2| + |b^\dagger_1 b^*_1| + |b^*_2 b^\dagger_2|)$.

    Define $\omega = \dist(A,B)$.
    Using the definitions of well separation and rules of a metric space we have:
    \[
        |a^*_1b^*_1| \geq \omega,\quad |b^*_2a^*_2| \geq \omega,\quad
        |a^*_1a^*_2|\leq \frac{\omega}{s},\quad |b^\dagger_1 b^*_1|\leq \frac{\omega}{s},\quad
        |b^*_2 b^\dagger_2|\leq \frac{\omega}{s},\quad |b^\dagger_1 b^\dagger_2|\geq 0,
    \]

    We can then conclude the following:
    \begin{align*}
        L - L' &= (|a^*_1 b^*_1| + |b^*_2 a^*_2| + |b^\dagger_1 b^\dagger_2|) - (|a^*_1 a^*_2| + |b^\dagger_1 b^*_1| + |b^*_2 b^\dagger_2|) \\
        &\geq (|a^*_1 b^*_1| + |b^*_2 a^*_2|) - (|a^*_1 a^*_2| + |b^\dagger_1 b^*_1| + |b^*_2 b^\dagger_2|) \\ 
        &\geq 2\omega - (|a^*_1 a^*_2| + |b^\dagger_1 b^*_1| + |b^*_2 b^\dagger_2|) \\
        &\geq 2\omega - \frac{3\omega}{s} \\
        &= \omega \left(2 - \frac{3}{s}\right)
    \end{align*}

    Thus if $s > \frac{3}{2}$ then we have that $L - L' > 0$ which is a contradiction since we assumed that the original paths were optimal (ie $L-L' \leq 0$).
    If $s = \frac{3}{2}$ then we have that $L - L' \geq 0$ and since by optimality we have that $L - L' \leq 0$ we can conclude that $L - L' = 0$ and therefore the newly constructed paths are also an optimal solution.
    Since this new paths removed the only two edges connecting $A$ to $B$ we can conclude that there exists an optimal solution where neither path has an edge connecting $A$ to $B$.

    Combining case 1 and case 2 we have that if $s \geq \frac{3}{2}$ then there exists an optimal solution where neither path has an edge connecting $A$ to $B$ and that if $s > \frac{3}{2}$ then all optimal solutions must satisfy this property.
    Furthermore it follows from case 1 that if $s \geq 1$ then there exists an optimal solution where there are at most two edges connecting $A$ to $B$ and if $s > 1$ then all optimal solutions must satisfy this property.
\end{proof}

\begin{thm}\label{thm:wsp_mpath_fixed_all_a}
    Suppose we have two sets of points $A$ and $B$ that are $s$-well-separated where $s \geq \frac{3}{2}$. 
    Then the multiple optimal hamiltonian paths problem of $k \geq 2$ paths where each paths endpoints are fixed and all endpoints are in $A$ has a solution where a single path visits all points in $B$ consecutively and all other paths only have points in $A$.
    Furthermore if $s > \frac{3}{2}$ then all optimal solutions must satisfy this property.
\end{thm}
\begin{proof} % gemini checked for major flaws
    Suppose that for some optimal solution there are at least two paths that have edges connecting $A$ to $B$.
    Then we can define two new sets of points $A', B'$ where $A'$ is the set of points in $A$ that are visited by these two paths and $B'$ is the set of points in $B$ that are visited by these two paths.
    Then we can note that $A'$ and $B'$ are still $s$-well-separated since they are subsets of $A$ and $B$ respectively.
    We can then define the endpoints of the first path as $a',a'' \in A$ and the endpoints of the second path as $a''',a'''' \in A$.
    If $s > \frac{3}{2}$ then this would contradict \Cref{lem:wsp_twopath_fixed_aa_aa} since it states that if $s > \frac{3}{2}$ then all optimal solutions must have only a single path with edges connecting $A$ to $B$.
    If $s = \frac{3}{2}$ then by \Cref{lem:wsp_twopath_fixed_aa_aa} there exists an optimal solution to the subproblem of the points in these two paths where only a single path has two edges connecting $A$ to $B$ and the other path has no edges connecting $A$ to $B$.
    We can then repeatedly apply this logic to keep generating optimal solutions with fewer paths connecting $A$ to $B$ until we have an optimal solution where only a single path has edges connecting $A$ to $B$ and all other paths have no edges connecting $A$ to $B$.

    Finally we can apply \Cref{thm:wsp_path_aa} to the path that has edges connecting $A$ to $B$ to construct a solution where this path visits all points in $B$ consecutively and furthermore if $s > \frac{3}{2}$ then all optimal solutions must satisfy this property.
\end{proof}

\begin{corollary}\label{cor:wsp_mtour_all_a}
    Suppose we have two sets of points $A$ and $B$ that are $s$-well-separated where $s \geq \frac{3}{2}$.
    Then the multiple optimal travelling salesman problem of $k \geq 2$ tours which has all starting and ending nodes in $A$ has a optimal solution where all tours but one stay within $A$ and the remaining tour visits all points in $B$ consecutively.
    Furthermore, if $s > \frac{3}{2}$ then all optimal solutions must satisfy this property.
\end{corollary}

\begin{thm}\label{thm:wsp_mpath_fixed_mixed}
    Suppose we have two sets of points $A$ and $B$ that are $s$-well-separated where $s \geq \frac{3}{2}$. 
    Then the multiple optimal hamiltonian paths problem of $k \geq 2$ paths where each paths endpoints are fixed and at least one path has an endpoint in $A$ and at least one path has an endpoint in $B$ has a solution where paths which have both an endpoint in $A$ and an endpoint in $B$ have a single edge connecting $A$ to $B$ and all other paths only have points in $A$ or only have points in $B$ (if they start/end in $A$ resp $B$).
    Furthermore if $s > \frac{3}{2}$ then all optimal solutions must satisfy this property.
\end{thm}
\begin{proof} % gemini checked for major flaws
    Assume that for some optimal solution there exists a path whose endpoints are just in $A$ xor $B$ and that path has an edge connecting $A$ to $B$.
    WLOG define the endpoints of the path to be $a',a'' \in A$ (otherwise just flip the names for $A$ and $B$).
    Choose another path which has at least on endpoint in $B$.
    We can then define $A'$ to be the set of points in $A$ that are visited by these two paths and $B'$ to be the set of points in $B$ that are visited by these two paths.
    Then we can note that $A'$ and $B'$ are still $s$-well-separated by \Cref{lem:wsp_subset}.

    \textbf{Case 1:} This path has both endpoints in $B$.
    Then we can apply \Cref{lem:wsp_twopath_fixed_aa_bb} to construct a solution where neither path has an edge connecting $A$ to $B$ and furthermore if $s > \frac{3}{2}$ then all optimal solutions must satisfy this property.

    \textbf{Case 2:} This path has one endpoint in $A$ and one endpoint in $B$.
    Then we can apply \Cref{lem:wsp_twopath_fixed_aa_ab} to construct a solution where only the second path has edges connecting $A$ to $B$ and furthermore if $s > \frac{3}{2}$ then all optimal solutions must satisfy this property.

    By repeatedly applying cases 1 and 2 we can construct a solution where paths that have both endpoints in $A$ xor $B$ never cross from $A$ to $B$. 
    Next take some path which has both one endpoint $a' \in A$ and one endpoint $b' \in B$ (if it exists).
    Define $A'$ to be the set of points in $A$ that are visited by this path and $B'$ to be the set of points in $B$ that are visited by this path.
    Then we can note that $A'$ and $B'$ are still $s$-well-separated by \Cref{lem:wsp_subset}.
    It then follows from \Cref{thm:wsp_path_ab} that the optimal solution must have that this path only has a single edge connecting $A$ to $B$.

    Therefore we have that if $s \geq \frac{3}{2}$ then there exists an optimal solution where paths which have both an endpoint in $A$ and an endpoint in $B$ have a single edge connecting $A$ to $B$ and all other paths only have points in $A$ or only have points in $B$ (if they start/end in $A$ resp $B$) and if $s > \frac{3}{2}$ then all optimal solutions must satisfy this property.
\end{proof}

\begin{corollary}\label{cor:wsp_mtour_mixed}
    Suppose we have two sets of points $A$ and $B$ that are $s$-well-separated where $s \geq \frac{3}{2}$.
    Then the multiple optimal travelling salesman problem of $k \geq 2$ paths which has a mixture of starting nodes in $A$ and $B$ has a solution where there are no edges connecting $A$ to $B$ in any of the tours.
    Furthermore if $s > \frac{3}{2}$ then all optimal solutions must satisfy this property.
\end{corollary}

\subsubsection{Single WSP Entry/Exit Set Multi Path Analysis}

\input{../MSThesis/theorems/wsp_twopath_aa_aa.tex}

\input{../MSThesis/theorems/wsp_twopath_aa_ab.tex}

\input{../MSThesis/theorems/wsp_twopath_aa_bb.tex}

\input{../MSThesis/theorems/wsp_twopath_ab_ab.tex}

\input{../MSThesis/theorems/wsp_twopath_ab_ab_tightness.tex}

\input{../MSThesis/theorems/wsp_mpath_start_a_end_a.tex}

\input{../MSThesis/theorems/wsp_mpath_start_end_mixed.tex}

\subsubsection{Single WSP Endpoint Set Multi Path Analysis}

\input{../MSThesis/theorems/wsp_twopath_aaaa.tex}

\input{../MSThesis/theorems/wsp_twopath_aaab.tex}

\input{../MSThesis/theorems/wsp_twopath_aabb.tex}

\begin{thm}\label{thm:wsp_mpath_all_a}
    Suppose we have two sets of points $A$ and $B$ that are $s$-well-separated where $s \geq 1$ and we have a multiple optimal hamiltonian paths problem with $k \geq 2$ paths where all endpoints are in $A$.
    Then there exists an optimal solution where a single path visits all points in $B$ consecutively and all other paths only have points in $A$.
    Furthermore if $s > 1$ then all optimal solutions must satisfy this property.
\end{thm}
\begin{proof}
    Consider some optimal solution to the problem.
    Suppose that there are at least two paths that have edges connecting $A$ to $B$.
    Then we can define two new sets of points $A', B'$ where $A'$ is the set of points in $A$ that are visited by these two paths and $B'$ is the set of points in $B$ that are visited by these two paths.
    Then we can note that $A'$ and $B'$ are still $s$-well-separated since they are subsets of $A$ and $B$ respectively.
    By \Cref{lem:wsp_twopath_aaaa} we know that there exists an optimal solution to the subproblem of the points in these two paths where only a single path has two edges connecting $A$ to $B$ and the other path has no edges connecting $A$ to $B$ and furthermore if $s > 1$ then all optimal solutions to this subproblem must satisfy this property.
    We can then apply this logic repeatedly to construct a solution where only a single path has edges connecting $A$ to $B$ and all other paths have no edges connecting $A$ to $B$.
    We can then apply \Cref{thm:wsp_path_aa} to the path that has edges connecting $A$ to $B$ to construct a solution where this path visits all points in $B$ consecutively and furthermore if $s > 1$ then all optimal solutions must satisfy this property.
\end{proof}

\begin{thm} 
    Suppose we have two sets of points $A$ and $B$ that are $s$-well-separated where $s \geq 1$ and we have a multiple optimal hamiltonian paths problem with $k \geq 2$ paths where there are an odd number of endpoints in $A$.
    Then there exists an optimal solution where a single path has exactly one edge connecting $A$ to $B$ and all other paths have no edges connecting $A$ to $B$.
    Furthermore if $s > 1$ then all optimal solutions must satisfy this property.
\end{thm}
\begin{proof}
    Consider some optimal solution to the problem.
    Suppose that there are at least two paths that have endpoints in different sets (one in $A$ and the other in $B$).
    Suppose that the first path has endpoints $a', b'$ and the second path has endpoints $a'', b''$.
    Then we can note that there must exist one edge of the form $(a^*, b^*)$ in the first path and one edge of the form $(a^\dagger, b^\dagger)$ in the second path where $a^*, a^\dagger \in A$ and $b^*, b^\dagger \in B$.
    This path has the following partial structure:
    $$a' \leadsto a^* \to b^* \leadsto b' \quad\quad a'' \leadsto a^\dagger \to b^\dagger \leadsto b''$$

    Since the paths are optimal we know that they satisfy the 2-optimality condition implying the following:
    $$|a^* b^*| + |a^\dagger b^\dagger| \leq |a^* a^\dagger| + |b^* b^\dagger|$$
    This is to say that the following paths must have a combined length at least as long as the original paths:
    $$a' \leadsto a^* \to a^\dagger \leadsto a'' \quad\quad b' \leadsto b^* \to b^\dagger \leadsto b''$$
    Since $A$ and $B$ are $s$-well-separated we know that $\forall i,j,k; |a_i a_j| \leq \frac{1}{s} |a_i b_k|$ (and vice versa for $B$). 
    From this we can conclude the following:
    \begin{align*}
        |a^* b^*| + |a^\dagger b^\dagger| &\leq |a^* a^\dagger| + |b^* b^\dagger| \\
        &\leq \frac{1}{s} |a^* b^*| + \frac{1}{s} |a^\dagger b^\dagger| \\
        &= \frac{1}{s} (|a^* b^*| + |a^\dagger b^\dagger|)
    \end{align*}
    This implies that $1 \leq \frac{1}{s}$ which in turn implies that $s \leq 1$.
    If $s > 1$ this is a contradiction and therefore any solution where there are at least two paths with endpoints in different sets cannot be optimal.
    If $s = 1$ then we can conclude that $|a^* b^*| + |a^\dagger b^\dagger| = |a^* a^\dagger| + |b^* b^\dagger|$ which implies that the newly constructed paths are also optimal.
    It follows that all but one path with endpoints in different sets can be replaced by paths with endpoints in the same set and thus there exists an optimal solution where only a single path has endpoints in different sets.

    We can then take this path with endpoints in different sets and any other path with endpoints in the same set and create two new sets of points $A', B'$ where $A'$ is the set of points in $A$ that are visited by these two paths and $B'$ is the set of points in $B$ that are visited by these two paths.
    Then we can note that $A'$ and $B'$ are still $s$-well-separated since they are subsets of $A$ and $B$ respectively.
    We can then apply \Cref{lem:wsp_twopath_aaab} to construct a solution to the subproblem of the points in these two paths where one has a single edge connecting $A$ to $B$ and the other has no edges connecting $A$ to $B$ and furthermore if $s > 1$ then all optimal solutions to this subproblem must satisfy this property.
    We can then apply this logic repeatedly to construct a solution where only a single path has an edge connecting $A$ to $B$ and all other paths have no edges connecting $A$ to $B$.
\end{proof}

\begin{thm}
    Suppose we have two sets of points $A$ and $B$ that are $s$-well-separated where $s \geq \frac{3}{2}$ and we have a multiple optimal hamiltonian paths problem with $k \geq 2$ paths where there are a non-zero even number of endpoints in $A$ and a non-zero even number of endpoints in $B$.
    Then there exists an optimal solution where there are no edges connecting $A$ to $B$ in any of the paths.
    Furthermore if $s > \frac{3}{2}$ then all optimal solutions must satisfy this property.
\end{thm}
\begin{proof}
    Consider some optimal solution to the problem.
    Suppose that there are at least two paths that have endpoints in different sets (one in $A$ and the other in $B$).
    Let the first path have endpoints $a', b'$ and the second path have endpoints $a'', b''$.
    Next we can define $A', B'$ where $A'$ is the set of points in $A$ that are visited by these two paths and $B'$ is the set of points in $B$ that are visited by these two paths.
    Then we can note that $A'$ and $B'$ are still $s$-well-separated since they are subsets of $A$ and $B$ respectively.
    We can then apply \Cref{lem:wsp_twopath_aabb} to construct a solution to the subproblem of the points in these two paths where each path has its endpoints in the same set furthermore if $s > \frac{3}{2}$ then all optimal solutions to this subproblem must satisfy this property.
    We can then apply this logic repeatedly to construct a solution where there are no paths with endpoints in different sets.
    
    Then for any path with endpoints both in $A$ (resp. $B$) we can take any other path with endpoints both in $B$ (resp. $A$) and apply \Cref{lem:wsp_twopath_aabb} to construct a solution to the subproblem of the points in these two paths where neither path has an edge connecting $A$ to $B$ and furthermore if $s > \frac{3}{2}$ then all optimal solutions to this subproblem must satisfy this property.
    We can then apply this logic repeatedly to construct a solution where there are no edges connecting $A$ to $B$ in any of the paths.
\end{proof}


\subsubsection{Binary Nested WSP Single Path Analysis}

\begin{lem} % I think this has to do with showing that one of these in on the "left" and the other is on the "right" of the current parent pair and thus any third pair would already be in one of these directions
    Suppose we have a given set of points $A$ which is individually $s$-well-separated from the sets of points $B,C,D$.
    
    If we were connected to two subseteq sets of $A$ direct pair it immediately follows that those cannot contribute.

    Suppose two were subseteq sets of $A$s parent pairs. [somehow this leads to a big problem] probably that those two themselves are Well separated from each other and thus we can 
\end{lem}

% TODO: Finish this proof by showing more than three sets is impossible
\begin{thm}\label{thm:bin_wspd_tour_structure}
    Let $P$ be a set of points and suppose we have a WSPD $\PairDecomp_s$ of $P$ with separation factor $s \geq 1$. 
    Then there must exist a optimal TSP tour over $P$ such that for any $k$-binary WSP $(A_i, B_i) \in \PairDecomp_s$ there is exactly zero, one, or two edges in the tour connecting $A_i,B_i$. Furthermore no more than four edges exit $A_i \cup B_i$.
    If $s > 1$ then all optimal TSP tours satisfy these properties.
\end{thm}
\begin{proof}
    % MARK: - The top level base case
    This can be proved via induction. 
    For the top level pair of $(A_t, B_t) \in \PairDecomp_s$ where $A_t \cup B_t = P$ the result follows directly from \Cref{cor:wsp_tsp} which implies there exists a tour where exactly two edges connect $A_t$ to $B_t$ and if $s > 1$ then all optimal tours have exactly two edges connecting $A_t$ to $B_t$.
    
    % MARK: - The second level base case
    Next consider the second level cluster $(A', B') \in \PairDecomp_s$. Let $C_t = A' \cup B'$ which is well separated from $D_t = P \setminus C_t$.
    By the previous case we know that there are exactly two edges exiting $C_t$ in some optimal TSP tour.
    WLOG let one of these edges be of the form $(a_*, d_*)$ where $a_* \in A'$ and $d_* \in D_t$ (we assume that $A'$ contains a connecting edge and if not then we can just swap $A'$ and $B'$).
    We proceed with a case analysis.

    \textbf{Case 1:} The second edge exiting $C_t$ is of the form $(b_*, d_\dagger)$ where $b_* \in B'$ and $d_\dagger \in D_t$.
    Since there are no other edges exiting $C_t$ we know that the optimal TSP tour must contain a hamiltonian path on $A_t$ with endpoints $a_*$ and $b_*$.
    By \Cref{thm:wsp_path_ab} we know that there exists an optimal path which has exactly one edge connecting $A'$ to $B'$ and that all optimal paths must satisfy this if $s > 1$.

    \textbf{Case 2:} The second edge exiting $C_t$ is of the form $(a_\dagger, d_\dagger)$ where $a_\dagger \in A'$ and $d_\dagger \in D_t$.
    Since there are no other edges exiting $C_t$ we know that the optimal TSP tour must contain a hamiltonian path on $A'$ with endpoints $a_*$ and $a_\dagger$.
    By \Cref{thm:wsp_path_aa} we know that there exists an optimal path which has exactly two edges connecting $A'$ to $B'$ and that all optimal paths must satisfy this if $s > 1$.

    % MARK: - The general case
    Now consider the third level or lower pair $(A_i, B_i) \in \PairDecomp_s$ of depth $i$ and let the parent pair to be $(C_p, D_p) \in \PairDecomp_s$ where $C_p = A_i \cup B_i$. 
    Assume by induction that there is some optimal tour where there is exactly zero, one, or two edges connecting $C_p$ to $D_p$ and furthermore let the number of edges exiting $C_p \cup D_p$ be $\ell \leq 2(i-2)$.

    It then immediately follows that the number of edges is at most $\ell + 2 \leq 2(i-1)$ since we can have at most two edges connecting $C_p$ to $D_p$ and at most two edges connecting $A_i$ to $B_i$.

    
\end{proof}


There may be a way to deterministically say which of the two cases will occur based on the structure of the WSPD but I have not yet found it.

\subsection{Moving beyond binary WSPDs}
\subsubsection{Triple Pairwise Separation}

\emph{TODO: Update these to have $s=1.0$}

\begin{thm}\label{thm:triple_wsp_path_aa}
    Let $P = A \cup B \cup C$ be a set of points where $A,B,C$ are pairwise $s$-well-separated with $s > 1$. 
    Then a hamiltonian path which starts on a point $a' \in A$ and ends on a point $a'' \in A$ must have only three or four edges connecting the three sets.
\end{thm}
\begin{proof}
    To prove this we can observe two different cases. 

    \textbf{Case 1:} The optimal hamiltonian path on $P$ contains at least one edge connecting each pair of sets.
    WLOG let these edges be $(a_1, b_1), (b_2, c_1), (c_2, a_2)$ ordered as they appear in the path (if $c$ occurs before $b$ then flip the names).
    Assume by contradiction that the optimal hamiltonian path from $a'$ to $a''$ over $P$ requires more than four edges connecting the three sets.
    Then there must either be another loop connecting the three sets or there must be at least two more edges connecting one of the pairs of sets.

    \textbf{Subcase 1a:} There is another loop connecting the three sets both in the same orientation.
    Then the optimal path must contain at least six edges connecting the three sets.
    WLOG let these edges be $(a_1, b_1)$, $(b_2, c_1)$, $(c_2, a_2)$, $(a_3, b_3)^*$, $(b_4, c_3)^*$, $(c_4, a_4)^*$ ordered as they appear in the path. This gives the following partial structure of the path:
    $$a' \to a_1 \to b_1 \leadsto b_2 \to c_1 \leadsto c_2 \to a_2 \leadsto a_3 \to b_3 \leadsto b_4 \to c_3 \leadsto c_4 \to a_4 \leadsto a''$$

    Since the path is optimal we know that it satisfies the 6-optimality condition implying the following:
    $|a_1 b_1| + |b_2 c_1| + |c_2 a_2| + |a_3 b_3| + |b_4 c_3| + |c_4 a_4| \leq |a_1 a_2| + |a_3 b_3| + |b_4 b_1| + |b_2 c_1| + |c_2 c_3| + |c_4 a_4|$.
    We can then rearrange this to get:
    \begin{align*}
        |a_1 b_1| + |c_2 a_2| + |b_4 c_3| &\leq |a_1 a_2| + |b_4 b_1| + |c_2 c_3| \\
        &\leq \frac{1}{s} |a_1 b_1| + \frac{1}{s} |c_2 a_2| + \frac{1}{s} |b_4 c_3| \\
        &= \frac{1}{s} (|a_1 b_1| + |c_2 a_2| + |b_4 c_3|)
    \end{align*}
    This implies that $1 \leq \frac{1}{s}$ which in turn implies that $s \leq 1$ contradicting our original statement that $s > 1$.

    \textbf{Subcase 1b:} There is another loop connecting the three sets in the opposite orientation.
    Then the optimal path must contain at least six edges connecting the three sets.
    WLOG let these edges be $(a_1, b_1)$, $(b_2, c_1)$, $(c_2, a_2)$, $(a_3, c_3)^*$, $(c_4, b_3)^*$, $(b_4, a_4)^*$ ordered as they appear in the path. This gives the following partial structure of the path:
    $$a' \to a_1 \to b_1 \leadsto b_2 \to c_1 \leadsto c_2 \to a_2 \leadsto a_3 \to c_3 \leadsto c_4 \to b_3 \leadsto b_4 \to a_4 \leadsto a''$$

    Since the path is optimal we know that it satisfies the 6-optimality condition implying the following:
    $|a_1 b_1| + |b_2 c_1| + |c_2 a_2| + |a_3 c_3| + |c_4 b_3| + |b_4 a_4| \leq |a_1 a_2| + |a_3 c_3| + |c_4 c_2| + |c_1 b_2| + |b_1 b_3| + |b_4 a_4|$.
    We can then rearrange this to get:
    \begin{align*}
        |a_1 b_1| + |c_2 a_2| + |c_4 b_3| &\leq |a_1 a_2| + |c_2 c_4| + |b_1 b_3| \\
        &\leq \frac{1}{s} |a_1 b_1| + \frac{1}{s} |c_2 a_2| + \frac{1}{s} |c_4 b_3| \\
        &= \frac{1}{s} (|a_1 b_1| + |c_2 a_2| + |c_4 b_3|)
    \end{align*}
    This implies that $1 \leq \frac{1}{s}$ which in turn implies that $s \leq 1$ contradicting our original statement that $s > 1$.

    \textbf{Subcase 1c:} There must be at least two additional edges at the end causing a loop from $A$ to $B$.
    Then the optimal hamiltonian path must contain at least five edges connecting the three sets.
    WLOG let these edges be $(a_1, b_1)$, $(b_2, c_1)$, $(c_2, a_2)$, $(a_3, b_3)^*$, $(b_4, a_4)^*$ ordered as they appear in the path. 
    This gives the following partial structure of the path:
    $$a' \to a_1 \to b_1 \leadsto b_2 \to c_1 \leadsto c_2 \to a_2 \leadsto a_3 \to b_3 \leadsto b_4 \to a_4 \leadsto a''$$

    Since the path is optimal we know that it satisfies the 2-optimality condition implying the following:
    \begin{align*}
        |a_1 b_1| + |a_3 b_3| &\leq |a_1 a_3| + |b_1 b_3| \\
        &\leq \frac{1}{s} |a_1 b_1| + \frac{1}{s} |a_3 b_3| \\
        &= \frac{1}{s} (|a_1 b_1| + |a_3 b_3|)
    \end{align*}
    This implies that $1 \leq \frac{1}{s}$ which in turn implies that $s \leq 1$ contradicting our original statement that $s > 1$.

    \textbf{Subcase 1d:} There must be at least two additional edges at the end causing a loop from $A$ to $C$.
    Then the optimal hamiltonian path must contain at least five edges connecting the three sets.
    WLOG let these edges be $(a_1, b_1)$, $(b_2, c_1)$, $(c_2, a_2)$, $(a_3, c_3)^*$, $(c_4, a_4)^*$ ordered as they appear in the path.
    This gives the following partial structure of the path:
    $$a' \to a_1 \to b_1 \leadsto b_2 \to c_1 \leadsto c_2 \to a_2 \leadsto a_3 \to c_3 \leadsto c_4 \to a_4 \leadsto a''$$

    Since the path is optimal we know that it satisfies the 2-optimality condition implying the following:
    \begin{align*}
        |c_2 a_2| + |c_4 a_4| &\leq |c_2 c_4| + |a_2 a_4| \\
        &\leq \frac{1}{s} |c_2 a_2| + \frac{1}{s} |c_4 a_4| \\
        &= \frac{1}{s} (|c_2 a_2| + |c_4 a_4|)
    \end{align*}
    This implies that $1 \leq \frac{1}{s}$ which in turn implies that $s \leq 1$ contradicting our original statement that $s > 1$.


    \textbf{Subcase 1e:} There must be at least two additional edges connecting $A,B$ before starting the primary loop.
    Then the optimal hamiltonian path must contain at least five edges connecting the three sets.
    WLOG let these edges be $(a_3, b_3)^*$, $(b_4, a_4)^*$, $(a_1, b_1)$, $(b_2, c_1)$, $(c_2, a_2)$ ordered as they appear in the path.
    This gives the following partial structure of the path:
    $$a' \to a_3 \to b_3 \leadsto b_4 \to a_4 \leadsto a_1 \to b_1 \leadsto b_2 \to c_1 \leadsto c_2 \to a_2 \leadsto a''$$

    Since the path is optimal we know that it satisfies the 2-optimality condition implying the following:
    \begin{align*}
        |a_3 b_3| + |a_1 b_1| &\leq |a_3 a_1| + |b_3 b_1| \\
        &\leq \frac{1}{s} |a_3 b_3| + \frac{1}{s} |a_1 b_1| \\
        &= \frac{1}{s} (|a_3 b_3| + |a_1 b_1|)
    \end{align*}
    This implies that $1 \leq \frac{1}{s}$ which in turn implies that $s \leq 1$ contradicting our original statement that $s > 1$.

    \textbf{Subcase 1f:} There must be at least two additional edges connecting $A,C$ before starting the primary loop.
    Then the optimal hamiltonian path must contain at least five edges connecting the three sets.
    WLOG let these edges be $(a_3, c_3)^*$, $(c_4, a_4)^*$, $(a_1, b_1)$, $(b_2, c_1)$, $(c_2, a_2)$ ordered as they appear in the path.
    This gives the following partial structure of the path:
    $$a' \to a_3 \to c_3 \leadsto c_4 \to a_4 \leadsto a_1 \to b_1 \leadsto b_2 \to c_1 \leadsto c_2 \to a_2 \leadsto a''$$

    Since the path is optimal we know that it satisfies the 2-optimality condition implying the following:
    \begin{align*}
        |c_4 a_4| + |c_2 a_2| &\leq |c_4 c_2| + |a_4 a_2| \\
        &\leq \frac{1}{s} |c_4 a_4| + \frac{1}{s} |c_2 a_2| \\
        &= \frac{1}{s} (|c_4 a_4| + |c_2 a_2|)
    \end{align*}
    This implies that $1 \leq \frac{1}{s}$ which in turn implies that $s \leq 1$ contradicting our original statement that $s > 1$.

    \textbf{Subcase 1g:} There is a loop back from $B$ to $C$ before the primary loop is ended.
    Then the optimal hamiltonian path must contain at least five edges connecting the three sets.
    WLOG let these edges be $(a_1, b_1)$, $(b_2, c_1)$, $(c_4, b_3)^*$, $(b_4, c_3)^*$, $(c_2, a_2)$ ordered as they appear in the path.
    This gives the following partial structure of the path:
    $$a' \to a_1 \to b_1 \leadsto b_2 \to c_1 \leadsto c_4 \to b_3 \leadsto b_4 \to c_3 \leadsto c_2 \to a_2 \leadsto a''$$

    Since the path is optimal we know that it satisfies the 2-optimality condition implying the following:
    \begin{align*}
        |b_2 c_1| + |b_4 c_3| &\leq |b_2 b_4| + |c_1 c_3| \\
        &\leq \frac{1}{s} |b_2 c_1| + \frac{1}{s} |b_4 c_3| \\
        &= \frac{1}{s} (|b_2 c_1| + |b_4 c_3|)
    \end{align*}
    This implies that $1 \leq \frac{1}{s}$ which in turn implies that $s \leq 1$ contradicting our original statement that $s > 1$.

    \textbf{Case 2:} The optimal TSP tour on $P$ does not contain an edge connecting each pair of sets.
    \textbf{Subcase 2a:} WLOG let there only edges connecting $A$ to $B$ and $B$ to $C$ ($C$ and $B$ can be alternated symmetrically).
    Then the optimal hamiltonian path must contain at least four edges connecting the three sets.
    WLOG let these edges be $(a_1, b_1)$, $(b_2, c_1)$, $(c_2, b_3)$, $(b_4, a_2)$ ordered as they appear in the path.

    \textbf{Subsubcase 2a.i:} There is an extra loop connecting $A$ to $B$ before ending the path.
    Then the optimal hamiltonian path must contain at least six edges connecting the three sets.
    WLOG let these edges be $(a_1, b_1)$, $(b_5, a_3)^*$, $(a_4, b_6)^*$, $(b_2, c_1)$, $(c_2, b_3)$, $(b_4, a_2)$ ordered as they appear in the path.
    This gives the following partial structure of the path:
    $$a' \leadsto a_1 \to b_1 \leadsto b_5 \to a_3 \leadsto a_4 \to b_6 \leadsto b_2 \to c_1 \leadsto c_2 \to b_3 \leadsto b_4 \to a_2 \leadsto a''$$

    Since the path is optimal we know that it satisfies the 2-optimality condition implying the following:
    \begin{align*}
        |a_1 b_1| + |a_4 b_6| &\leq |a_1 a_4| + |b_1 b_6| \\
        &\leq \frac{1}{s} |a_1 b_1| + \frac{1}{s} |a_4 b_6| \\
        &= \frac{1}{s} (|a_1 b_1| + |a_4 b_6|)
    \end{align*}
    This implies that $1 \leq \frac{1}{s}$ which in turn implies that $s \leq 1$ contradicting our original statement that $s > 1$.
    
    \textbf{Subsubcase 2a.ii:} There is an extra loop connecting $B$ to $C$ before ending the path.
    Then the optimal hamiltonian path must contain at least six edges connecting the three sets.
    WLOG let these edges be $(a_1, b_1)$, $(b_2, c_1)$, $(c_3, b_5)^*$, $(b_6, c_4)^*$, $(c_2, b_3)$, $(b_4, a_2)$ ordered as they appear in the path.
    This gives the following partial structure of the path:
    $$a' \leadsto a_1 \to b_1 \leadsto b_2 \to c_1 \leadsto c_3 \to b_5 \leadsto b_6 \to c_4 \leadsto c_2 \to b_3 \leadsto b_4 \to a_2 \leadsto a''$$

    Since the path is optimal we know that it satisfies the 2-optimality condition implying the following:
    \begin{align*}
        |b_2 c_1| + |b_6 c_4| &\leq |b_2 b_6| + |c_1 c_4| \\
        &\leq \frac{1}{s} |b_2 c_1| + \frac{1}{s} |b_6 c_4| \\
        &= \frac{1}{s} (|b_2 c_1| + |b_6 c_4|)
    \end{align*}
    This implies that $1 \leq \frac{1}{s}$ which in turn implies that $s \leq 1$ contradicting our original statement that $s > 1$.

    \textbf{Subsubcase 2a.iii:} There is an extra loop connecting $A,B$ after ending the path.
    Then the optimal hamiltonian path must contain at least six edges connecting the three sets.
    WLOG let these edges be $(a_1, b_1)$, $(b_2, c_1)$, $(c_2, b_3)$, $(b_4, a_2)$, $(a_3, b_5)^*$, $(b_6, a_4)^*$ ordered as they appear in the path.
    This gives the following partial structure of the path:
    $$a' \leadsto a_1 \to b_1 \leadsto b_2 \to c_1 \leadsto c_2 \to b_3 \leadsto b_4 \to a_2 \leadsto a_3 \to b_5 \leadsto b_6 \to a_4 \leadsto a''$$

    Since the path is optimal we know that it satisfies the 2-optimality condition implying the following:
    \begin{align*}
        |a_1 b_1| + |a_3 b_5| &\leq |a_1 a_3| + |b_1 b_5| \\
        &\leq \frac{1}{s} |a_1 b_1| + \frac{1}{s} |a_3 b_5| \\
        &= \frac{1}{s} (|a_1 b_1| + |a_3 b_5|)
    \end{align*}
    This implies that $1 \leq \frac{1}{s}$ which in turn implies that $s \leq 1$ contradicting our original statement that $s > 1$.

    \textbf{Subcase 2b:} Let there only be edges connecting $A$ to $B$ and $A$ to $C$.
    WLOG let these edges be $(a_1, b_1)$, $(b_2, a_2)$, $(a_3, c_1)$, $(c_2, a_4)$ ordered as they appear in the path (if $C$ occurs before $B$ then flip the names).

    \textbf{Subsubcase 2b.i:} There is an extra loop connecting $A$ to $B$ before ending the path.
    Then the optimal hamiltonian path must contain at least six edges connecting the three sets.
    WLOG let these edges be $(a_1, b_1)$, $(b_5, a_3)^*$, $(a_4, b_6)^*$, $(b_2, a_2)$, $(a_5, c_1)$, $(c_2, a_6)$ ordered as they appear in the path.
    This gives the following partial structure of the path:
    $$a' \leadsto a_1 \to b_1 \leadsto b_5 \to a_3 \leadsto a_4 \to b_6 \leadsto b_2 \to a_2 \leadsto a_5 \to c_1 \leadsto c_2 \to a_6 \leadsto a''$$

    Since the path is optimal we know that it satisfies the 2-optimality condition implying the following:
    \begin{align*}
        |a_1 b_1| + |a_4 b_6| &\leq |a_1 a_4| + |b_1 b_6| \\
        &\leq \frac{1}{s} |a_1 b_1| + \frac{1}{s} |a_4 b_6| \\
        &= \frac{1}{s} (|a_1 b_1| + |a_4 b_6|)
    \end{align*}
    This implies that $1 \leq \frac{1}{s}$ which in turn implies that $s \leq 1$ contradicting our original statement that $s > 1$.

    \textbf{Subsubcase 2b.ii:} There is an extra loop connecting $A,C$ after ending the path.
    Then the optimal hamiltonian path must contain at least six edges connecting the three sets.
    WLOG let these edges be $(a_1, b_1)$, $(b_2, a_2)$, $(a_3, c_1)$, $(c_2, a_4)$, $(a_5, c_3)^*$, $(c_4, a_6)^*$ ordered as they appear in the path.
    This gives the following partial structure of the path:
    $$a' \leadsto a_1 \to b_1 \leadsto b_2 \to a_2 \leadsto a_3 \to c_1 \leadsto c_2 \to a_4 \leadsto a_5 \to c_3 \leadsto c_4 \to a_6 \leadsto a''$$

    Since the path is optimal we know that it satisfies the 2-optimality condition implying the following:
    \begin{align*}
        |a_3 c_1| + |a_5 c_3| &\leq |a_3 a_5| + |c_1 c_3| \\
        &\leq \frac{1}{s} |a_3 c_1| + \frac{1}{s} |a_5 c_3| \\
        &= \frac{1}{s} (|a_3 c_1| + |a_5 c_3|)
    \end{align*}
    This implies that $1 \leq \frac{1}{s}$ which in turn implies that $s \leq 1$ contradicting our original statement that $s > 1$.

    \textbf{Subsubcase 2b.iii:} There is an extra loop connecting $A,B$ before ending the path.
    Then the optimal hamiltonian path must contain at least six edges connecting the three sets.
    WLOG let these edges be $(a_1, b_1)$, $(b_2, a_2)$, $(a_3, c_1)$, $(c_2, a_4)$, $(a_5, b_5)^*$, $(b_6, a_6)^*$ ordered as they appear in the path.
    This gives the following partial structure of the path:
    $$a' \leadsto a_1 \to b_1 \leadsto b_2 \to a_2 \leadsto a_3 \to c_1 \leadsto c_2 \to a_4 \leadsto a_5 \to b_5 \leadsto b_6 \to a_6 \leadsto a''$$

    Since the path is optimal we know that it satisfies the 2-optimality condition implying the following:
    \begin{align*}
        |a_1 b_1| + |a_5 b_5| &\leq |a_1 a_5| + |b_1 b_5| \\
        &\leq \frac{1}{s} |a_1 b_1| + \frac{1}{s} |a_5 b_5| \\
        &= \frac{1}{s} (|a_1 b_1| + |a_5 b_5|)
    \end{align*}
    This implies that $1 \leq \frac{1}{s}$ which in turn implies that $s \leq 1$ contradicting our original statement that $s > 1$.

    Therefore in all cases we have reached a contradiction when assuming that there are more than four edges connecting the three sets.
    Thus the optimal hamiltonian path from $a'$ to $a''$ over $P$ must have only three or four edges connecting the three sets.
\end{proof}

\begin{corollary}
    Let $P = A \cup B \cup C$ be a set of points where $A,B,C$ are pairwise $s$-well-separated with $s > 1$. 
    Then $A$, $B$, and $C$ must be TSP-separated.
\end{corollary}

\subsubsection{Generalization to \texorpdfstring{$k$}{k} pairwise separation}
\begin{thm}
    Let $P = \bigcup_{i=1}^n A_i$ be a set of points where $A_1, A_2, \ldots, A_n$ are all pairwise $s$-well-separated with $s \geq 1$ (and $n \geq 2$).
    Then there exists a optimal hamiltonian path over $P$ which starts at $a' \in A_i$ and ends at $a'' \in A_j$ where no set is ever directly connected to another set more than twice in the same direction.
    % Having no more than three edges connecting is a special case since >= 3 edges implies two in the same direction
    If $s > 1$ then all optimal hamiltonian paths must satisfy this property.
    % TODO: what if we got A to B to A then to C to B to C to A? we would have four paths connecting A to B
\end{thm}
\begin{proof}
    Consider some optimal hamiltonian path over $P$ which starts at $a' \in A_i$ and ends at $a'' \in A_j$.
    Assume there is two pairs of sets $B = A_k$ and $C = A_l$ which are directly connected by two edges in the same direction (where $k \neq l$).
    WLOG let these edges be $(b_1, c_1)$ and $(b_2, c_2)$ where $b_1, b_2 \in B$ and $c_1, c_2 \in C$ ordered as they appear in the path (if $C$ occurs before $B$ then flip the names).

    This gives the following partial structure of the path:
    $$a' \leadsto b_1 \to c_1 \leadsto b_2 \to c_2 \leadsto a''$$
    Since the path is optimal we know that it satisfies the 2-optimality condition implying that $|b_1 c_1| + |b_2 c_2| \leq |b_1 b_2| + |c_1 c_2|$.
    This is to say that the following path must be at least as long as the original path:
    $$a' \leadsto b_1 \to b_2 \leadsto c_1 \to c_2 \leadsto a''$$
    Since $C$ and $B$ are $s$-well-separated we know that $\forall i,j,k; |b_i b_j| \leq \frac{1}{s} |b_i c_k|$ (and vice versa for $C$). 
    From this we can conclude the following:
    \begin{align*}
        |b_1 c_1| + |b_2 c_2| &\leq |b_1 b_2| + |c_1 c_2| \\
        &\leq \frac{1}{s} |b_1 c_1| + \frac{1}{s} |b_2 c_2| \\
        &= \frac{1}{s} (|b_1 c_1| + |b_2 c_2|)
    \end{align*}
    This implies that $1 \leq \frac{1}{s}$ which in turn implies that $s \leq 1$.
    Thus if $s > 1$ then we have a contradiction and therefore no optimal hamiltonian path can have more than two edges directly connecting any pair of sets in the same direction.
    If $s = 1$ then we can conclude that $|b_1 c_1| + |b_2 c_2| = |b_1 b_2| + |c_1 c_2|$ and therefore the second path is also an optimal hamiltonian path over $P$.
    We can repeat this process over all pairs until every set is connected by no more than one edge to another set in a given direction.

    Therefore there exists an optimal hamiltonian path over $P$ which starts at $a' \in A_i$ and ends at $a'' \in A_j$ where no set is ever directly connected to another set more than twice in the same direction.
\end{proof}

% TODO: THIS PROOF IN NOT FINISHED
\begin{thm}
    Let $P = \bigcup_i A_i$ be a set of points where $A_1, A_2, \ldots, A_n$ are all pairwise $s$-well-separated with $s \geq 1$ (and $n \geq 2$).
    Then there exists a optimal hamiltonian path over $P$ which starts at $a' \in A_i$ and ends at $a'' \in A_j$ where no two non-overlapping subpaths of the hamiltonian path both start in the same set $A_k$ and end in the same set $A_l$ for any $k,l$ (where $k \neq l$).
    If $s > 1$ then all optimal hamiltonian paths must satisfy this property.
\end{thm}
\begin{proof}
    Consider some optimal hamiltonian path over $P$ which starts at $a' \in A_i$ and ends at $a'' \in A_j$.
    Assume there is two pairs of sets $B = A_k$ and $C = A_l$ which are directly connected by two edges in the same direction (where $k \neq l$).
    WLOG let these edges be $(b_1, c_1)$ and $(b_2, c_2)$ where $b_1, b_2 \in B$ and $c_1, c_2 \in C$ ordered as they appear in the path (if $C$ occurs before $B$ then flip the names).

    This can be shown via induction on the number of sets $n$.

    \textbf{Base case:} $n=2$. This follows directly from \Cref{thm:wsp_path_ab,thm:wsp_path_aa}.

    \textbf{Base case:} $n=3$. This follows directly from \Cref{thm:triple_wsp_path_aa,thm:triple_wsp_path_ab}.

    \textbf{Inductive step:} Assume the statement is true for all $n' < n$.

    Consider some optimal hamiltonian path over $P$ which starts at $a' \in A_i$ and ends at $a'' \in A_j$.
    Assume there is two pairs of sets $B = A_k$ and $C = A_l$ such that there are two non-overlapping subpaths of the hamiltonian path which both start in $B$ and end in $C$.
    Let one path start at $b_1$ and end at $c_1$ and the other path start at $b_2$ and end at $c_2$.

    Let $t$ represent the number of sets from $A_1, A_2, \ldots, A_n$ that are visited by the path from $b_1$ to $c_2$ (including $B$ and $C$).
    
    \textbf{Case 1:} $t < n$. Then we can then create a new instance $P'$ by only including the set of points visited by the path from $b_1$ to $c_2$.
    If $s > 1$ then it follows from the inductive hypothesis that this path could not be optimal.
    If $s = 1$ then it follows from the inductive hypothesis that there exists a separate optimal path that does not have two non-overlapping subpaths between $b_1$ and $c_2$.

    \textbf{Case 2:} $t = n$. Then the path from $b_1$ to $c_2$ visits all sets in $P$.
    They cannot both be direct connects as this follows from \Cref{thm:original}
\end{proof}

\section{Algorithms}

\begin{algorithm}[H]
\caption{\textsc{BestBalancedBiPartition}$(D,s)$}
\label{alg:balanced-bipartition}
\begin{algorithmic}[1]
\Require A metric distance matrix $D \in \mathbb{R}^{n \times n}$ (symmetric, triangle inequality), separation factor $s > 1$
\Ensure A bi-partition $(A,B)$ that is $s$-well-separated and maximally balanced, or \textsc{NULL} if no such partition exists

\State $E \gets \textsc{BuildMST}(D)$ \Comment{$n-1$ edges}
\State $\text{bestScore} \gets +\infty$
\State $\text{bestPartition} \gets \textsc{NULL}$

\ForAll{edges $e = (u,v)$ in $E$}
    \State $(A,B) \gets \textsc{ConnectedComponents}(E \setminus \{e\})$
    \State $\text{balance} \gets \max(|A|, |B|)$ \Comment{Balance score: size of larger cluster}
    \If{$\text{balance} \geq \text{bestScore}$}
        \State \textbf{continue} \Comment{Skip if not better than current best}
    \EndIf

    \State $\delta \gets w(e)$ \Comment{Single-linkage separation distance}

    \State $\operatorname{diam}(A) \gets \textsc{ComputeDiameter}(A,D)$
    \State $\operatorname{diam}(B) \gets \textsc{ComputeDiameter}(B,D)$
    \State $M \gets \max\{\operatorname{diam}(A), \operatorname{diam}(B)\}$

    \If{$\delta \ge s \cdot M$} \Comment{$s$-well-separated condition}
        \State $\text{bestScore} \gets \text{balance}$
        \State $\text{bestPartition} \gets (A,B)$
        \If{$\text{bestScore} = \lceil n/2 \rceil$} \Comment{Perfectly balanced partition found}
            \State \textbf{break}
        \EndIf
    \EndIf
\EndFor

\State \Return $\text{bestPartition}$
\end{algorithmic}
\end{algorithm}

Analyzing the program there are $n-1$ edges in the MST which takes $O(n^2)$ time to compute, and for each edge we compute the connected components in $O(n)$ time, compute the diameter of each component in $O(n^2)$ time, and check the separation condition in $O(1)$ time.
Thus the overall runtime is $O(n^3)$.

\begin{lem}\label[lem]{lem:mst-contains-min-edge}
    Let $(A,B)$ be a bi-partition of a set of points $P$ within a metric space.
    Then any MST of $P$ must contain at least on edge connecting the two clusters $A$ and $B$ that is minimal among all edges connecting $A$ and $B$.
\end{lem}
\begin{proof}
    Let $e = (u,v)$ be an edge representing the minimum distance between $A$ and $B$ where $u \in A$ and $v \in B$.
    Suppose that $e$ is not in some MST of $P$. 
    Then by adding $e$ to the MST we create a cycle.
    Since $u \in A$ and $v \in B$ there must exist another edge $e' = (u', v')$ in the cycle where $u' \in A$ and $v' \in B$.
    Since $e$ represents the minimum distance between $A$ and $B$ we have that $\omega(e) \leq \omega(e')$.
    However by the minimax property of the MST we know that $\omega(e') \leq \omega(e)$.
    Thus $\omega(e) = \omega(e')$ and therefore the MST contains an edge representing the minimum distance between $A$ and $B$.
\end{proof}

\begin{lem}\label[lem]{lem:multiple-cuts-bound-separation}
    Let $(A,B)$ be a bi-partition of a set of points $P$ within a metric space. Suppose that more than one cut of any MST is required to separate $A$ and $B$. Then $(A,B)$ cannot be $s$-well-separated for any $s > 1$.
\end{lem}
\begin{proof}
    Let $e = (u, v)$ be the edge in the MST representing the minimum distance between $A$ and $B$ where $u \in A$ and $v \in B$ which must exist by \Cref{lem:mst-contains-min-edge}.
    Since more than one cut is required to separate $A$ and $B$, there must be another edge $e' = (v', u')$ in the MST where $u' \in A$ and $v' \in B$.
    Let $M = \max\{\diam(A), \diam(B)\}$.
    Define $e^* = (u, u')$ then we know that $\omega(e^*) \leq M$ by the definition of diameter.
    However by the minimax property of the MST we have that $\omega(e) \leq \omega(e^*)$.
    Thus $\omega(e) \leq M$ and therefore $(A,B)$ cannot be $s$-well-separated for any $s > 1$.
\end{proof}

\begin{thm}
    If there exists a bi-partition of the points that is $s$-well-separated (where $s>1$), then \Cref{alg:balanced-bipartition} will find a bi-partition that is $s$-well-separated and maximally balanced.
\end{thm}
\begin{proof}
    It follows from \Cref{lem:multiple-cuts-bound-separation} that since $s>1$ then any bi-partition that is $s$-well-separated must be separable by a single cut of any MST.
    Since \Cref{alg:balanced-bipartition} iterates through all edges of the MST and checks the bi-partition induced by cutting that edge, it will find all bi-partitions that are $s$-well-separated.
    Among these bi-partitions, the algorithm keeps track of the one with the best balance score (size of larger cluster) and returns it at the end.
    Therefore if there exists a bi-partition that is $s$-well-separated, the algorithm will find and return a bi-partition that is $s$-well-separated and maximally balanced.
\end{proof}

\begin{thm}
    Any algorithm that finds the best balanced bi-partition that is $s$-well-separated for $s\leq1$ is NP-hard.
\end{thm}
\begin{proof}
    This can be shown by a reduction from the Partition problem.
    Given a multiset of $n$ positive integers $S = \{x_1, x_2, \dots, x_n\}$, the Partition problem asks whether $S$ can be
    partitioned into two subsets $S_1$ and $S_2$ such that the sum of the elements in each subset is equal.
    We can construct a special set of points $P$ in the following way: for each integer $x_i$ in $S$, create $x_i$ points that are clumped on a unique vertex of the $n$ simplex. 
    Assume that the solution to the partition problem has each partition made up of at least two integers (if not we can trivially check for this case).
    Since every clump lies on a unique vertex of the simplex, the diameter of any partition is 1 and the distance between any two partitions is 1 thus making the separation factor exactly 1.
    Therefore any bi-partition of $P$ is 1-well-separated.
    It can then be noted that after solving for the best balanced bi-partition of $P$ that is $s$-well-separated for $s \leq 1$, if the resulting bi-partition has both clusters with equal size, then the answer to the Partition problem is yes and if not then the answer must be no.
\end{proof}

\subsubsection{Second Version}

\begin{algorithm}[H]
\caption{\textsc{BalancedMetricSplit2}$(D,s,k)$}
\label{alg:balanced_metric_split2}
\begin{algorithmic}[1]
\Require A metric distance matrix $D \in \mathbb{R}^{n\times n}$ (symmetric, triangle inequality), separation factor $s>0$, and number of clusters $k \ge 2$
\Ensure A no more than $k$-partition $\mathcal{P}=\{S_1,\dots,S_k\}$ that maximizes balance such that every pair of clusters is $s$-well-separated (or a single cluster of all points is returned)

\State $\mathcal{M} \gets \{D_{ij} \mid 1 \le i < j \le n\}$ \Comment{candidate diameter thresholds}
\State $\mathcal{S}^\star \gets \big\{ \{1,2,\dots,n\} \big\}$; $\text{bestScore} \gets n$

\ForAll{$M \in \mathcal{M}$}
    \State $E_M \gets \{(i,j)\mid 1\le i<j\le n,\; D_{ij} < s\cdot M\}$
    \State $G_M \gets (V=\{1,\dots,n\}, E_M)$
    \State $\mathcal{C} \gets \textsc{ConnectedComponents}(G_M)$

    \If{$1 < |\mathcal{C}| \leq l$}
        \State \textbf{continue}
    \EndIf

    \Comment{Diameter feasibility: ensure each component has diameter $\le M$}
    \State $\text{ok} \gets \textbf{true}$
    \ForAll{$C \in \mathcal{C}$}
        \If{$\diam(C) > M$}
            \State $\text{ok} \gets \textbf{false}$
            \State \textbf{break}
        \EndIf
    \EndFor
    \If{\textbf{not} ok}
        \State \textbf{continue}
    \EndIf

    \Comment{Separation feasibility: ensure pairwise separation}
    \State $\text{sep} \gets \textbf{true}$
    \ForAll{distinct $C, C' \in \mathcal{C}$}
        \State $\Delta(C,C') \gets \min\{D_{xy}\mid x\in C,\; y\in C'\}$
        \If{$\Delta(C,C') < s \cdot \max(\diam(C),\diam(C'))$}
            \State $\text{sep} \gets \textbf{false}$
            \State \textbf{break}
        \EndIf
    \EndFor
    \If{\textbf{not} sep}
        \State \textbf{continue}
    \EndIf

    \Comment{Balance objective: minimize size of largest cluster}
    \State $\text{score} \gets \max_{C\in\mathcal{C}} \bigl|C\bigr|$
    \If{$\text{score} < \text{bestScore}$}
        \State $\text{bestScore} \gets \text{score}$
        \State $\mathcal{S}^\star \gets \mathcal{C}$
    \EndIf
\EndFor

\State \Return $\mathcal{S}^\star$
\end{algorithmic}
\end{algorithm}

\appendix
\section{Outdated proofs}
% this section is for if a better bound is found (ie proof is still valid but not as strong as desired)
\begin{thm}[There is a better bound]
    Suppose we have two sets of points $A$ and $B$ that are $s$-well-separated where $s > 2.5$. Then $A$ and $B$ must be TSP-separated.
\end{thm}

\begin{rem*}
    The previous corollary cannot be guaranteed for $s \leq 0.5$
\end{rem*}
\begin{proof}
    If $s \leq 0.5$ a evil problem can be created in 1d. 
    Consider $A=\{0,1\}$ and $B=\{0.5,1.5\}$. 
    The diameter of each set is 1 and distance between the two sets is 0.5 making them 0.5-well-separated.
    However the optimal TSP tour on $A \cup B$ is $0 \to 0.5 \to 1 \to 1.5 \to 0$ which would contradict \Cref{cor:wsp_tsp} if it applied for $s \leq 0.5$.
\end{proof}

\section{Dead ends}

$$\binom{n^2}{2} \cdot 2 \binom{(n/2)^2}{2} \cdot 4 \binom{(n/4)^2}{2} \cdot \dots$$
$$x^{\log(x)}$$

% these theorems are false

%Let $A$ and $B$ be two finite sets of points that are TSP-separated. 
%Let $\Cover_k(S)$ denote the optimal cost of $k$ disjoint hamiltonian paths covering all points in set $S$.
%Let $\BigLength_{2k}(A,B)$ be the sum of the $2k$ shortest edges between $A$ and $B$ without point replacement.
%Then, $\Cover_1(A) + \Cover_1(B) + \BigLength_{2}(A,B) \leq \Cover_k(A) + \Cover_k(B) + \BigLength_{2k}(A,B)$.

%Let $A$ and $B$ be two finite sets of points that are TSP-separated. Then $\dist(A,B) \geq \max(\lambda(A), \lambda(B))$.
%[Where $\lambda(A) = \max(|a_1 a_2|, |a_2 a_3|, \dots, |a_{n-1} a_n|)$ so that it represents the longest edge in the hamiltonian path through $A$ which is derived as a sutwopath of the optimal TSP tour on $A \cup B$. The same definition applies to $\lambda(B)$.]


% This proof didnt work since the edges can grow exponentially
%\begin{proof}
%    % MARK: - The top level case
%    This can be proved via induction. 
%    For the top level pair of $(A_t, B_t) \in \PairDecomp_s$ where $A_t \cup B_t = P$ the result follows directly from \Cref{cor:wsp_tsp} which implies there exists a tour where exactly two edges connect $A_t$ to $B_t$ and if $s > 1$ then all optimal tours have exactly two edges connecting $A_t$ to $B_t$.
    
%    % MARK: - The second level case
%    Next consider the second level cluster $(A', B') \in \PairDecomp_s$. Let $C_t = A' \cup B'$ which is well separated from $D_t = P \setminus C_t$.
%    By the previous case we know that there are exactly two edges exiting $C_t$ in some optimal TSP tour.
%    WLOG let one of these edges be of the form $(a_*, d_*)$ where $a_* \in A'$ and $d_* \in D_t$ (we assume that $A'$ contains a connecting edge and if not then we can just swap $A'$ and $B'$).
%    We proceed with a case analysis.

%    \textbf{Case 1:} The second edge exiting $C_t$ is of the form $(b_*, d_\dagger)$ where $b_* \in B'$ and $d_\dagger \in D_t$.
%    Since there are no other edges exiting $C_t$ we know that the optimal TSP tour must contain a hamiltonian path on $A_t$ with endpoints $a_*$ and $b_*$.
%    By \Cref{thm:wsp_path_ab} we know that there exists an optimal path which has exactly one edge connecting $A'$ to $B'$ and that all optimal paths must satisfy this if $s > 1$.

%    \textbf{Case 2:} The second edge exiting $C_t$ is of the form $(a_\dagger, d_\dagger)$ where $a_\dagger \in A'$ and $d_\dagger \in D_t$.
%    Since there are no other edges exiting $C_t$ we know that the optimal TSP tour must contain a hamiltonian path on $A'$ with endpoints $a_*$ and $a_\dagger$.
%    By \Cref{thm:wsp_path_aa} we know that there exists an optimal path which has exactly two edges connecting $A'$ to $B'$ and that all optimal paths must satisfy this if $s > 1$.

%    % MARK: - The nth level cas

%    % contributed
%    % 1 edge connect and 2 edges exiting
%    % 2 edges connect and 2 edges exiting

%    Now consider the third level or lower pair $(A_i, B_i) \in \PairDecomp_s$ and let the parent pair to be $(C_p, D_p) \in \PairDecomp_s$ where $C_p = A_i \cup B_i$. 
%    Assume by induction that there is some non empty subset of optimal tours where there is exactly zero, one, or two edges connecting $C_p$ to $D_p$ and correspondingly there are exactly four, two, or two edges exiting $C_p \cup D_p$ respectively.

%    %\textbf{Case 1:} There are zero edges connecting $C_p$ to $D_p$ in the optimal TSP tour and there are four edges exiting $C_p \cup D_p$.
%    %Then since each of $C_p$ and $D_p$ must be visited we know that there must be exactly two edges exiting $C_p$.
%    %WLOG let one of these edges be of the form $(a_*, e_*)$ where $a_* \in A_i$ and $e_* \notin C_p \cup D_p$.

%    %\textbf{Case 1a:} The second edge exiting $C_p$ is of the form $(b_*, e_\dagger)$ where $b_* \in B_i$ and $e_\dagger \notin C_p \cup D_p$.
%    %Since there are no other edges exiting $C_p$ we know that the optimal TSP tour must contain a hamiltonian path on $C_p$ with endpoints $a_*$ and $b_*$.
%    %By \Cref{thm:wsp_path_ab} we know that there exists an optimal path which has exactly one edge connecting $A_i$ to $B_i$ and that all optimal paths must satisfy this if $s > 1$.

%    %\textbf{Case 1b:} The second edge exiting $C_p$ is of the form $(a_\dagger, e_\dagger)$ where $a_\dagger \in A_i$ and $e_\dagger \notin C_p \cup D_p$.
%    %Since there are no other edges exiting $C_p$ we know that the optimal TSP tour must contain a hamiltonian path on $C_p$ with endpoints $a_*$ and $a_\dagger$.
%    %By \Cref{thm:wsp_path_aa} we know that there exists an optimal path which has exactly two edges connecting $A_i$ to $B_i$ and that all optimal paths must satisfy this if $s > 1$.

%    \textbf{Case 1:} There is exactly one edge connecting $C_p$ to $D_p$ and there are two edges exiting $C_p \cup D_p$.
%    Then there must be exactly two edges exiting $C_p$.
%    WLOG let one of these edges be of the form $(a_*, d_*)$ where $a_* \in A_i$ and $d_* \in D_p$ (the case where the edge is of the form $(b_*, d_*)$ is symmetric).

%    \textbf{Case 1a:} The second edge exiting $C_p$ is of the form $(b_*, e_*)$ where $b_* \in B_i$ and $e_* \notin C_p \cup D_p$.
%    Since there are no other edges exiting $C_p$ we know that the tour must contain a hamiltonian path on $C_p$ with endpoints $a_*$ and $b_*$.
%    By \Cref{thm:wsp_path_ab} we know that there exists an optimal path which has exactly one edge connecting $A_i$ to $B_i$ and that all optimal paths must satisfy this if $s > 1$.

%    \textbf{Case 1b:} The second edge exiting $C_p$ is of the form $(a_\dagger, e_*)$ where $a_\dagger \in A_i$ and $e_* \notin C_p \cup D_p$.
%    Since there are no other edges exiting $C_p$ we know that the tour must contain a hamiltonian path on $C_p$ with endpoints $a_*$ and $a_\dagger$.
%    By \Cref{thm:wsp_path_aa} we know that there exists an optimal path which has exactly two edges connecting $A_i$ to $B_i$ and that all optimal paths must satisfy this if $s > 1$.

%    % contributed
%    % 1 edge connect and 2 edges exiting 
%    % 2 edges connect and 2 edges exiting

%    \textbf{Case 2:} There are two edges connecting $C_p$ to $D_p$ in the optimal TSP tour and there are two edges exiting $C_p \cup D_p$.
    
%    % ae bd be bd
%    % ae ad be bd
%    % ae ae bd bd
%    % ae ae ad bd
%    % ae ad ad be
%    % ae ae ad ad

%    \textbf{Case 2a:} Both exit edges are of the form $(c_*, e_*)$, $(c_\dagger, e_\dagger)$ where $c_*, c_\dagger \in C_p$ and $e_*, e_\dagger \notin C_p \cup D_p$.
%    Then there are four edges which exit $C_p$.
%    WLOG assume that $A_i$ contains an edge of the form $(a_1, e_*)$ where $a_* \in A_i$ and $e_* \notin C_p \cup D_p$ (the case where $B_i$ contains an edge of this form is symmetric).

%    \textbf{Case 2a.i:} The other three edges are of the following form $(b_1, d_*)$, $(b_2, e_\dagger)$, $(b_3, d_\dagger)$ where $b_1, b_2, b_3 \in B_i$, $d_*, d_\dagger \in D_p$ and $e_\dagger \notin C_p \cup D_p$.
%    Since there are no other edges exiting $C_p$ we know that the optimal TSP tour must contain two hamiltonian paths on $C_p$ with endpoints $a_1,b_1,b_2,b_3$.
%    Then by \Cref{lem:wsp_twopath_aaab} we know that there exists an optimal combination of paths over $C_p$ with endpoints $a_1, b_1, b_2, b_3$ that has exactly one edge connecting $A_i$ to $B_i$ and that all optimal combinations of paths must satisfy this if $s > 1$.

%    \textbf{Case 2a.ii:} The other three edges are of the following form $(a_2, d_*)$, $(b_1, e_\dagger)$, $(b_2, d_\dagger)$ where $a_2 \in A_i$, $b_1, b_2 \in B_i$, $d_*, d_\dagger \in D_p$ and $e_\dagger \notin C_p \cup D_p$.
%    Since there are no other edges exiting $C_p$ we know that the optimal TSP tour must contain two hamiltonian paths on $C_p$ with endpoints $a_1, a_2, b_1, b_2$.
%    Then by \Cref{lem:wsp_twopath_aabb} we know that there exists an optimal combination of paths over $C_p$ with endpoints $a_1, a_2, b_1, b_2$ that has no edges connecting $A_i$ to $B_i$ and that all optimal combinations of paths must satisfy this if $s > \frac{3}{2}$.

%    \textbf{Case 2a.iii:} The other three edges are of the following form $(a_2, e_\dagger)$, $(b_1, d_*)$, $(b_2, d_\dagger)$ where $a_2 \in A_i$, $b_1, b_2 \in B_i$, $d_*, d_\dagger \in D_p$ and $e_\dagger \notin C_p \cup D_p$.
%    Since there are no other edges exiting $C_p$ we know that the optimal TSP tour must contain two hamiltonian paths on $C_p$ with endpoints $a_1, a_2, b_1, b_2$.
%    Then by \Cref{lem:wsp_twopath_aabb} we know that there exists an optimal combination of paths over $C_p$ with endpoints $a_1, a_2, b_1, b_2$ that has no edges connecting $A_i$ to $B_i$ and that all optimal combinations of paths must satisfy this if $s > \frac{3}{2}$.

%    \textbf{Case 2a.iv:} The other three edges are of the following form $(a_2, d_*)$, $(a_3, e_\dagger)$, $(b_1, d_\dagger)$ where $a_2, a_3 \in A_i$, $b_1 \in B_i$, $d_*, d_\dagger \in D_p$ and $e_\dagger \notin C_p \cup D_p$.
%    Since there are no other edges exiting $C_p$ we know that the optimal TSP tour must contain two hamiltonian paths on $C_p$ with endpoints $a_1, a_2, a_3, b_1$.
%    Then by \Cref{lem:wsp_twopath_aaab} we know that there exists an optimal combination of paths over $C_p$ with endpoints $a_1, a_2, a_3, b_1$ that has exactly one edge connecting $A_i$ to $B_i$ and that all optimal combinations of paths must satisfy this if $s > 1$.

%    \textbf{Case 2a.v:} The other three edges are of the following form $(a_2, d_*)$, $(a_3, d_\dagger)$, $(b_1, e_\dagger)$ where $a_2, a_3 \in A_i$, $b_1 \in B_i$, $d_*, d_\dagger \in D_p$ and $e_\dagger \notin C_p \cup D_p$.
%    Since there are no other edges exiting $C_p$ we know that the optimal TSP tour must contain two hamiltonian paths on $C_p$ with endpoints $a_1, a_2, a_3, b_1$.
%    Then by \Cref{lem:wsp_twopath_aaab} we know that there exists an optimal combination of paths over $C_p$ with endpoints $a_1, a_2, a_3, b_1$ that has exactly one edge connecting $A_i$ to $B_i$ and that all optimal combinations of paths must satisfy this if $s > 1$.

%    \textbf{Case 2a.vi:} The other three edges are of the following form $(a_2, d_*)$, $(a_3, d_\dagger)$, $(a_4, e_\dagger)$ where $a_2, a_3, a_4 \in A_i$, $d_*, d_\dagger \in D_p$ and $e_\dagger \notin C_p \cup D_p$.
%    Since there are no other edges exiting $C_p$ we know that the optimal TSP tour must contain two hamiltonian paths on $C_p$ with endpoints $a_1, a_2, a_3, a_4$.
%    Then by \Cref{lem:wsp_twopath_aaaa} we know that there exists an optimal combination of paths over $C_p$ with endpoints $a_1, a_2, a_3, a_4$ that has exactly two edges connecting $A_i$ to $B_i$ and that all optimal combinations of paths must satisfy this if $s > 1$.

%    % contributed
%    % 1 edge connect and one has 3 exits and the other has 1 exit
%    % 0 edge connect and each have 2 exits
%    % 0 edge connect and each have 2 exits
%    % 1 edge connect and one has 3 exits and the other has 1 exit
%    % 1 edge connect and one has 3 exits and the other has 1 exit
%    % 2 edge connect and one has 4 exits and the other has 0 exits

%    \textbf{Case 2b:} Both exit edges are of the form $(d_*, e_*)$, $(d_\dagger, e_\dagger)$ where $d_*, d_\dagger \in D_p$ and $e_*, e_\dagger \notin C_p \cup D_p$.
%    We define the edges connecting $C_p$ to $D_p$ to be $(c', d')$ and $(c'', d'')$ where $c', c'' \in C_p$ and $d', d'' \in D_p$.
%    Since there are no other edges exiting $C_p$ we know that the optimal TSP tour must contain a single hamiltonian path on $C_p$ with endpoints $c', c''$.
    
%    \textbf{Case 2b.i:} One of $c', c''$ is in $A_i$ and the other is in $B_i$.
%    By \Cref{thm:wsp_path_ab} we know that there exists an optimal path which has exactly one edge connecting $A_i$ to $B_i$ and that all optimal paths must satisfy this if $s > 1$.

%    \textbf{Case 2b.ii:} Both of $c', c''$ are in $A_i$ or both of $c', c''$ are in $B_i$.
%    By \Cref{thm:wsp_path_aa} we know that there exists an optimal path which has exactly two edges connecting $A_i$ to $B_i$ and that all optimal paths must satisfy this if $s > 1$.

%    % contributed
%    % 1 edge connect and 2 edges exiting (fixed)
%    % 2 edge connect and 2 edges exiting (fixed)

%    \textbf{Case 2c:} One of the exit edges is of the form $(c_*, e_*)$ where $c_* \in C_p$ and $e_* \notin C_p \cup D_p$ and the other edge is of the form $(d_\dagger, e_\dagger)$ where $d_\dagger \in D_p$ and $e_\dagger \notin C_p \cup D_p$.
%    This is not possible since such a situation would require an odd number of edges connecting $C_p$ to $D_p$ which is a contradiction since we are in the case where there are two edges connecting $C_p$ to $D_p$.

%    \textbf{Case 3:} There are zero edges connecting $C_p$ to $D_p$ in the optimal TSP tour and each set has two edges exiting $C_p \cup D_p$.
%    Then there must be exactly two edges exiting $C_p$.
%    Let these edges be of the form $(c_*, e_*)$, $(c_\dagger, e_\dagger)$ where $c_*, c_\dagger \in C_p$ and $e_*, e_\dagger \notin C_p \cup D_p$.
%    Since there are no other edges exiting $C_p$ we know that the optimal TSP tour must contain a hamiltonian path on $C_p$ with endpoints $c_*$ and $c_\dagger$.

%    \textbf{Case 3a:} One of $c_*, c_\dagger$ is in $A_i$ and the other is in $B_i$.
%    By \Cref{thm:wsp_path_ab} we know that there exists an optimal path which has exactly one edge connecting $A_i$ to $B_i$ and that all optimal paths must satisfy this if $s > 1$.

%    \textbf{Case 3b:} Both of $c_*, c_\dagger$ are in $A_i$ or both of $c_*, c_\dagger$ are in $B_i$.
%    By \Cref{thm:wsp_path_aa} we know that there exists an optimal path which has exactly two edges connecting $A_i$ to $B_i$ and that
%    all optimal paths must satisfy this if $s > 1$.

%    % contributed
%    % 1 edge connect and 2 edges exiting (fixed)
%    % 2 edge connect and 2 edges exiting (fixed)

%    \textbf{Case 4:} There is one edge connecting $C_p$ to $D_p$ in the optimal TSP tour and there are four edges exiting $C_p \cup D_p$ where three of these edges come from one of $C_p$ or $D_p$ and the other edge comes from the other set.
%    Let the connecting edge be of the form $(c_*, d_*)$ where $c_* \in C_p$ and $d_* \in D_p$.

%    \textbf{Case 4a:} There is only a single exit edge from $C_p$.
%    Let this edge be of the form $(c', e')$ where $c' \in C_p$ and $e' \notin C_p \cup D_p$.
%    Since there are no other edges exiting $C_p$ we know that the optimal TSP tour must contain a hamiltonian path on $C_p$ with endpoints $c_*$ and $c'$.

%    \textbf{Case 4a.i:} One of $c_*, c'$ is in $A_i$ and the other is in $B_i$.
%    By \Cref{thm:wsp_path_ab} we know that there exists an optimal path which has exactly one edge connecting $A_i$ to $B_i$ and that all optimal paths must satisfy this if $s > 1$.

%    \textbf{Case 4a.ii:} Both of $c_*, c'$ are in $A_i$ or both of $c_*, c'$ are in $B_i$.
%    By \Cref{thm:wsp_path_aa} we know that there exists an optimal path which has exactly two edges connecting $A_i$ to $B_i$ and that all optimal paths must satisfy this if $s > 1$.

%    % contributed
%    % 1 edge connect and 2 edges exiting (fixed)
%    % 2 edge connect and 2 edges exiting (fixed)

%    \textbf{Case 4b:} There are three exit edges from $C_p$.
%    Let these edges be of the form $(c', e')$, $(c'', e'')$, $(c''', e''')$ where $c', c'', c''' \in C_p$ and $e', e'', e''' \notin C_p \cup D_p$.
%    Since there are no other edges exiting $C_p$ we know that the optimal TSP tour must contain two hamiltonian paths on $C_p$ with endpoints $c_*, c', c'', c'''$.

%    \textbf{Case 4b.i:} All of $c', c'', c''', c_*$ are all exclusively in $A_i$ or $B_i$.
%    Then by \Cref{lem:wsp_twopath_aaaa} we know that there exists an optimal combination of paths over $C_p$ with endpoints $c_*, c', c'', c'''$ that has exactly two edges connecting $A_i$ to $B_i$ and that all optimal combinations of paths must satisfy this if $s > 1$.

%    \textbf{Case 4b.ii:} Exactly one of $c', c'', c''', c_*$ is in $A_i$ (resp. $B_i$) and the other three are in $B_i$ (resp. $A_i$).
%    Then by \Cref{lem:wsp_twopath_aaab} we know that there exists an optimal combination of paths over $C_p$ with endpoints $c_*, c', c'', c'''$ that has exactly one edge connecting $A_i$ to $B_i$ and that all optimal combinations of paths must satisfy this if $s > 1$.

%    \textbf{Case 4b.iii:} Exactly two of $c', c'', c''', c_*$ are in $A_i$ and the other two are in $B_i$.
%    Then by \Cref{lem:wsp_twopath_aabb} we know that there exists an optimal combination of paths over $C_p$ with endpoints $c_*, c', c'', c'''$ that has no edges connecting $A_i$ to $B_i$ and that all optimal combinations of paths must satisfy this if $s > \frac{3}{2}$.

%    % contributed
%    % 2 edges connect and one has 4 exits and the other has 0 exits
%    % 1 edge connects and one has 3 exits and the other has 1 exit
%    % 0 edge connect and each have 2 exits

%    \textbf{Case 5:} There are two edges connecting $C_p$ to $D_p$ in the optimal TSP tour and there are four edges exiting $C_p \cup D_p$ where all four edges come from one of $C_p$ or $D_p$.
%    Let the two connecting edges be of the form $(c_*, d_*)$, $(c_\dagger, d_\dagger)$ where $c_*, c_\dagger \in C_p$ and $d_*, d_\dagger \in D_p$.

%    \textbf{Case 5a:} None of the exit edges come from $C_p$.
%    Since there are no other edges exiting $C_p$ we know that the optimal TSP tour must contain a hamiltonian path on $C_p$ with endpoints $c_*$ and $c_\dagger$.

%    \textbf{Case 5a.i:} One of $c_*, c_\dagger$ is in $A_i$ and the other is in $B_i$.
%    By \Cref{thm:wsp_path_ab} we know that there exists an optimal path which has exactly one edge connecting $A_i$ to $B_i$ and that all optimal paths must satisfy this if $s > 1$.

%    \textbf{Case 5a.ii:} Both of $c_*, c_\dagger$ are in $A_i$ or both of $c_*, c_\dagger$ are in $B_i$.
%    By \Cref{thm:wsp_path_aa} we know that there exists an optimal path which has exactly two edges connecting $A_i$ to $B_i$ and that all optimal paths must satisfy this if $s > 1$.

%    % contributed
%    % 1 edge connect and two edges exiting (fixed)
%    % 2 edge connect and two edges exiting (fixed)

%    \textbf{Case 5b:} Two of the exit edges come from $C_p$.
%    Let these edges be of the form $(c_1, e_1)$, $(c_2, e_2)$ where $c_1, c_2 \in C_p$ and $e_1, e_2 \notin C_p \cup D_p$.
%    Since there are no other edges exiting $C_p$ we know that the optimal TSP tour must contain a hamiltonian path on $C_p$ with endpoints $c_*, c_\dagger, c_1, c_2$.

%    \textbf{Case 5b.i:} All of $c_*, c_\dagger, c_1, c_2$ are all exclusively in $A_i$ or $B_i$.
%    Then by \Cref{lem:wsp_twopath_aaaa} we know that there exists an optimal combination of paths over $C_p$ with endpoints $c_*, c_\dagger, c_1, c_2$ that has exactly two edges connecting $A_i$ to $B_i$ and that all optimal combinations of paths must satisfy this if $s > 1$.

%    \textbf{Case 5b.ii:} Exactly one of $c_*, c_\dagger, c_1, c_2$ is in $A_i$ (resp. $B_i$) and the other three are in $B_i$ (resp. $A_i$).
%    Then by \Cref{lem:wsp_twopath_aaab} we know that there exists an optimal combination of paths over $C_p$ with endpoints $c_*, c_\dagger, c_1, c_2$ that has exactly one edge connecting $A_i$ to $B_i$ and that all optimal combinations of paths must satisfy this if $s > 1$.

%    \textbf{Case 5b.iii:} Exactly two of $c_*, c_\dagger, c_1, c_2$ are in $A_i$ and the other two are in $B_i$.
%    Then by \Cref{lem:wsp_twopath_aabb} we know that there exists an optimal combination of paths over $C_p$ with endpoints $c_*, c_\dagger, c_1, c_2$ that has no edges connecting $A_i$ to $B_i$ and that all optimal combinations of paths must satisfy this if $s > \frac{3}{2}$.

%    % contributed
%    % 2 edges connect and one has 4 exits and the other has 0 exits
%    % 1 edge connect and one has 3 exits and the other has 1 exit
%    % 0 edge connect and each have 2 exits

%    \textbf{Case 5c:} Four of the exit edges come from $C_p$.
%    Let these edges be of the form $(c_1, e_1)$, $(c_2, e_2)$, $(c_3, e_3)$, $(c_4, e_4)$ where $c_1, c_2, c_3, c_4 \in C_p$ and $e_1, e_2, e_3, e_4 \notin C_p \cup D_p$.
%    Since there are no other edges exiting $C_p$ we know that the optimal TSP tour must contain three hamiltonian paths on $C_p$ with endpoints $c_*, c_\dagger, c_1, c_2, c_3, c_4$.




%    %It then follows that the following edges must exist $(a_*, e_*)$, $(a_\dagger, d_*)$, $(b_*, e_\dagger)$, $(b_\dagger, d_\dagger)$ where $a_*, a_\dagger \in A_i$, $b_*, b_\dagger \in B_i$, $d_*, d_\dagger \in D_p$, and $e_*, e_\dagger \notin C_p \cup D_p$.

%    %Assume by contradiction that no optimal hamiltonian path with endpoints $c_*,c_\dagger \in C_p$ over $C_p \cup D_p$ visits all of the points in $A_i$ and $B_i$ consecutively.
%    %Then there must be at least two edges connecting $A_i$ to $B_i$.



%    %Let $(c_*, d_*)$ be an edge connecting $C_p$ to $D_p$ in the optimal TSP tour.
%    %WLOG assume that $c_* \in A_i$ and define $a_* = c_*$ (the case where $c_* \in B_i$ is symmetric).
%    %We proceed with a case analysis.

%    %\textbf{Case 1:} There are two edges connecting $C_p$ to $D_p$ in the optimal TSP tour and the second edge is of the form $(b_*, d_\dagger)$ where $b_* \in B_i$ and $d_\dagger \in D_p$.
%    %Since there are no other edges exiting $C_p$ we know that the optimal TSP tour must contain a hamiltonian path on $C_p$ with endpoints $a_*$ and $b_*$.
%    %By Theorem~\ref{thm:wsp_path_ab} we know that this path must visit all points in $A_i$ before any point in $B_i$ (or vice versa).
%    %Thus there must be exactly one edge connecting $A_i$ to $B_i$. And exactly two edges $(a_*, d_*)$, $(b_*, d_\dagger)$ exiting $C_p$ as desired.

%    %\textbf{Case 2:} There are two edges connecting $C_p$ to $D_p$ in the optimal TSP tour and the second edge is of the form $(a_\dagger, d_\dagger)$ where $a_\dagger \in A_i$ and $d_\dagger \in D_p$.
%    %Since there are no other edges exiting $C_p$ we know that the optimal TSP tour must contain a hamiltonian path on $C_p$ with endpoints $a_*$ and $a_\dagger$.
%    %By Theorem~\ref{thm:wsp_path_aa} we know that this path must contain exactly two edges between $A_i$ and $B_i$. Furthermore there are exactly two edges $(a_*, d_*)$, $(a_\dagger, d_\dagger)$ exiting $C_p$ as desired.

%    %\textbf{Case 3:} There is only one edge connecting $C_p$ to $D_p$ in the optimal TSP tour and the other edge is of the form $(a_\dagger, e_*)$ where $a_\dagger \in A_i$ and $e_* \notin C_p \cup D_p$.
%    %Since there are no other edges exiting $C_p$ we know that the optimal TSP tour must contain a hamiltonian path on $C_p$ with endpoints $a_*$ and $a_\dagger$.
%    %By Theorem~\ref{thm:wsp_path_aa} we know that this path must only contain two edges between $A_i$ and $B_i$. Furthermore there are exactly two edges $(a_*, d_*)$, $(a_\dagger, e_*)$ exiting $C_p$ as desired.

%    %\textbf{Case 4:} There is only one edge connecting $C_p$ to $D_p$ in the optimal TSP tour and the other edge is of the form $(b_*, e_*)$ where $b_* \in B_i$ and $e_* \notin C_p \cup D_p$.
%    %Since there are no other edges exiting $C_p$ we know that the optimal TSP tour must contain a hamiltonian path on $C_p$ with endpoints $a_*$ and $b_*$.
%    %By Theorem~\ref{thm:wsp_path_ab} we know that this path must visit all points in $A_i$ before any point in $B_i$ (or vice versa).
%    %Thus there must be exactly one edges connecting $A_i$ to $B_i$. And exactly two edges $(a_*, d_*)$, $(b_*, e_*)$ exiting $C_p$ as desired.

%    Therefore by induction we have shown that for any given $k$-binary pair $(A_i, B_i) \in \PairDecomp_s$ there must exist exactly one or two edges in the optimal TSP tour on $P$ connecting them. Furthermore, there are exactly two edges exiting $A_i \cup B_i$ if they are not the root clusters.
%\end{proof}

\end{document}